% Options for packages loaded elsewhere
\PassOptionsToPackage{unicode}{hyperref}
\PassOptionsToPackage{hyphens}{url}
\documentclass[
]{article}
\usepackage{xcolor}
\usepackage{amsmath,amssymb}
\setcounter{secnumdepth}{-\maxdimen} % remove section numbering
\usepackage{iftex}
\ifPDFTeX
  \usepackage[T1]{fontenc}
  \usepackage[utf8]{inputenc}
  \usepackage{textcomp} % provide euro and other symbols
\else % if luatex or xetex
  \usepackage{unicode-math} % this also loads fontspec
  \defaultfontfeatures{Scale=MatchLowercase}
  \defaultfontfeatures[\rmfamily]{Ligatures=TeX,Scale=1}
\fi
\usepackage{lmodern}
\ifPDFTeX\else
  % xetex/luatex font selection
\fi
% Use upquote if available, for straight quotes in verbatim environments
\IfFileExists{upquote.sty}{\usepackage{upquote}}{}
\IfFileExists{microtype.sty}{% use microtype if available
  \usepackage[]{microtype}
  \UseMicrotypeSet[protrusion]{basicmath} % disable protrusion for tt fonts
}{}
\makeatletter
\@ifundefined{KOMAClassName}{% if non-KOMA class
  \IfFileExists{parskip.sty}{%
    \usepackage{parskip}
  }{% else
    \setlength{\parindent}{0pt}
    \setlength{\parskip}{6pt plus 2pt minus 1pt}}
}{% if KOMA class
  \KOMAoptions{parskip=half}}
\makeatother
\usepackage{color}
\usepackage{fancyvrb}
\newcommand{\VerbBar}{|}
\newcommand{\VERB}{\Verb[commandchars=\\\{\}]}
\DefineVerbatimEnvironment{Highlighting}{Verbatim}{commandchars=\\\{\}}
% Add ',fontsize=\small' for more characters per line
\newenvironment{Shaded}{}{}
\newcommand{\AlertTok}[1]{\textcolor[rgb]{1.00,0.00,0.00}{\textbf{#1}}}
\newcommand{\AnnotationTok}[1]{\textcolor[rgb]{0.38,0.63,0.69}{\textbf{\textit{#1}}}}
\newcommand{\AttributeTok}[1]{\textcolor[rgb]{0.49,0.56,0.16}{#1}}
\newcommand{\BaseNTok}[1]{\textcolor[rgb]{0.25,0.63,0.44}{#1}}
\newcommand{\BuiltInTok}[1]{\textcolor[rgb]{0.00,0.50,0.00}{#1}}
\newcommand{\CharTok}[1]{\textcolor[rgb]{0.25,0.44,0.63}{#1}}
\newcommand{\CommentTok}[1]{\textcolor[rgb]{0.38,0.63,0.69}{\textit{#1}}}
\newcommand{\CommentVarTok}[1]{\textcolor[rgb]{0.38,0.63,0.69}{\textbf{\textit{#1}}}}
\newcommand{\ConstantTok}[1]{\textcolor[rgb]{0.53,0.00,0.00}{#1}}
\newcommand{\ControlFlowTok}[1]{\textcolor[rgb]{0.00,0.44,0.13}{\textbf{#1}}}
\newcommand{\DataTypeTok}[1]{\textcolor[rgb]{0.56,0.13,0.00}{#1}}
\newcommand{\DecValTok}[1]{\textcolor[rgb]{0.25,0.63,0.44}{#1}}
\newcommand{\DocumentationTok}[1]{\textcolor[rgb]{0.73,0.13,0.13}{\textit{#1}}}
\newcommand{\ErrorTok}[1]{\textcolor[rgb]{1.00,0.00,0.00}{\textbf{#1}}}
\newcommand{\ExtensionTok}[1]{#1}
\newcommand{\FloatTok}[1]{\textcolor[rgb]{0.25,0.63,0.44}{#1}}
\newcommand{\FunctionTok}[1]{\textcolor[rgb]{0.02,0.16,0.49}{#1}}
\newcommand{\ImportTok}[1]{\textcolor[rgb]{0.00,0.50,0.00}{\textbf{#1}}}
\newcommand{\InformationTok}[1]{\textcolor[rgb]{0.38,0.63,0.69}{\textbf{\textit{#1}}}}
\newcommand{\KeywordTok}[1]{\textcolor[rgb]{0.00,0.44,0.13}{\textbf{#1}}}
\newcommand{\NormalTok}[1]{#1}
\newcommand{\OperatorTok}[1]{\textcolor[rgb]{0.40,0.40,0.40}{#1}}
\newcommand{\OtherTok}[1]{\textcolor[rgb]{0.00,0.44,0.13}{#1}}
\newcommand{\PreprocessorTok}[1]{\textcolor[rgb]{0.74,0.48,0.00}{#1}}
\newcommand{\RegionMarkerTok}[1]{#1}
\newcommand{\SpecialCharTok}[1]{\textcolor[rgb]{0.25,0.44,0.63}{#1}}
\newcommand{\SpecialStringTok}[1]{\textcolor[rgb]{0.73,0.40,0.53}{#1}}
\newcommand{\StringTok}[1]{\textcolor[rgb]{0.25,0.44,0.63}{#1}}
\newcommand{\VariableTok}[1]{\textcolor[rgb]{0.10,0.09,0.49}{#1}}
\newcommand{\VerbatimStringTok}[1]{\textcolor[rgb]{0.25,0.44,0.63}{#1}}
\newcommand{\WarningTok}[1]{\textcolor[rgb]{0.38,0.63,0.69}{\textbf{\textit{#1}}}}
\usepackage{longtable,booktabs,array}
\newcounter{none} % for unnumbered tables
\usepackage{calc} % for calculating minipage widths
% Correct order of tables after \paragraph or \subparagraph
\usepackage{etoolbox}
\makeatletter
\patchcmd\longtable{\par}{\if@noskipsec\mbox{}\fi\par}{}{}
\makeatother
% Allow footnotes in longtable head/foot
\IfFileExists{footnotehyper.sty}{\usepackage{footnotehyper}}{\usepackage{footnote}}
\makesavenoteenv{longtable}
\setlength{\emergencystretch}{3em} % prevent overfull lines
\providecommand{\tightlist}{%
  \setlength{\itemsep}{0pt}\setlength{\parskip}{0pt}}
\usepackage{bookmark}
\IfFileExists{xurl.sty}{\usepackage{xurl}}{} % add URL line breaks if available
\urlstyle{same}
\hypersetup{
  pdftitle={Semiotic Taps: Lightweight Adapter Modules for Bifurcation Detection in Frozen Language Models},
  pdfauthor={James Burton Lancaster},
  hidelinks,
  pdfcreator={LaTeX via pandoc}}

\title{Semiotic Taps: Lightweight Adapter Modules for Bifurcation
Detection in Frozen Language Models}
\author{James Burton Lancaster}
\date{April 27, 2026}

\begin{document}
\maketitle

\subsection{Abstract}\label{abstract}

Large language models trained on web-scale corpora absorb the semiotic
bifurcations embedded in their data. These are divergent interpretant
chains in which the same sign carries incompatible meanings across
discourse communities. Such models have no mechanism to detect,
represent, or respond to this divergence. We introduce the
Semiotic-Reflexive Transformer Adapter (SRT-Adapter), a lightweight
architecture (\textasciitilde14.5M trainable parameters, 0.19\% of a 7B
backbone) that bolts semiotic awareness onto any frozen causal language
model without modifying its embeddings, attention, or output head. The
adapter operates through four modules that \emph{tap} hidden states at
selected backbone layers: (1) a \textbf{Community Discovery Head} that
performs unsupervised soft clustering of discourse communities from
early-layer representations; (2) \textbf{Metapragmatic Attention Heads}
(MAH) that compute divergence vectors quantifying where meaning forks
under community-conditioned interpretation; (3) a \textbf{Reflexive
Recurrent Module} (RRM) that tracks accumulated semiotic divergence
through a per-position GRU meta-state and optionally injects small
corrections into the backbone stream via FiLM modulation; and (4) a
\textbf{Bifurcation Estimation Network} (BEN) that estimates a
continuous reflexivity coefficient \(\hat{r}\) and a binary semiotic
regime (subcritical/supercritical) at each token position. Grounded in
Peircean semiotics and the pitchfork bifurcation model of political
polarization (Lancaster, 2025), the architecture treats the frozen
backbone as a substrate on which semiotic processes are an emergent,
measurable phenomenon. Training combines the backbone's native
cross-entropy with auxiliary losses on chain-of-interpretants
prediction, bifurcation regression, regime classification, divergence
health, community entropy, and supervised-contrastive separation on both
the community channel (v5) and the metapragmatic divergence channel
(v6), together with ListNet ranking on \(\hat{r}\) and a chain-residual
auxiliary floor. The corpus is 1M Reddit samples spanning 35 discourse
communities with per-token reflexivity annotations. We report a
five-generation empirical arc on a Qwen 2.5-7B backbone. v5 establishes
the basic capability set on a five-probe suite: cross-entropy
preservation (CE = 2.63 vs.~unadapted 2.71), unsupervised community
retrieval (recall@1 = 0.36, \(12.6\times\) random on a 35-class task),
counterfactual community decoding (zero disagreement on factual prompts,
0.95 mean disagreement on contested topics), zero-shot hallucination
signal on TruthfulQA (mean \(\hat{r}\) AUROC = 0.573), and regime
calibration (ECE = \(9\times10^{-4}\), AUROC = 0.99 on 351K tokens;
replicated on v8a at ECE \(= 9.1\times 10^{-4}\), §5.5). The headline
finding is in v8a: removing the 32-prototype mixing layer entirely,
leaving the 64-D encoder output as the community vector, leaves CE
unchanged (Δ = +0.0001 nats) while raising Reddit recall@1 from 0.413 to
0.484, raising archetype recall@1 on an out-of-distribution 33-class
taxonomy from 0.149 to 0.230 (\(7.6\times\) chance), nearly doubling the
within/between cosine ratio (1.006 → 2.016), and expanding trajectory
anisotropy by \(\sim\)\(325\times\). v8b falsifies the
``sharper-supcon'' hypothesis on this architecture: doubling the
contrastive weight and halving the temperature partially undoes v8a's
gains. The encoder, not the prototype basis, was doing the
discriminative work. Two post-release results sharpen the architectural
reading. A cross-backbone raw-hidden probe (§5.12) shows the
discourse-community signal v8a amplifies is latent in
\texttt{Qwen/Qwen3-8B} and \texttt{mistralai/Mistral-7B-v0.3} at
comparable strength to Qwen 2.5-7B, indicating the adapter targets a
substrate that is not Qwen-specific. An interiority study on the v8a
checkpoint (§5.13) shows the inject channels function as a fixed-size
BOS register whose amplitude is stable across an \(8\times\) length
sweep (peak-to-peak \(\leq 4.8\%\)), while the regime classifier reads
content-borne signal off non-BOS tokens and remains calibrated.

\textbf{Keywords:} semiotic adapter, bifurcation detection,
metapragmatic attention, interpretant chains, reflexive recurrence,
discourse community discovery, frozen backbone, pitchfork bifurcation,
Peircean semiotics

\begin{center}\rule{0.5\linewidth}{0.5pt}\end{center}

\subsection{1. Introduction}\label{introduction}

\subsubsection{1.1 The Problem}\label{the-problem}

Language models are semiotic infrastructure. Their outputs enter
interpretant chains alongside human-authored signs, shaping subsequent
interpretation in ways neither users nor developers can fully trace. Yet
the training paradigm that produces these systems is semiotically naive:
it optimizes for the conditional probability of the next token, a
surface-level objective that captures co-occurrence patterns while
remaining structurally blind to the interpretive processes that make
those patterns meaningful.

The consequence is that when a language model encounters a contested
sign such as ``freedom,'' ``justice,'' or ``woke,'' it produces text
that is fluent within a particular attractor basin without representing
the fact that the sign indexes opposed interpretive communities. The
model does not know it is in a bifurcation zone. It cannot tell you.

Current alignment methods (RLHF, DPO, Constitutional AI) intervene
downstream, constraining outputs after the model has already
internalized a bifurcated semiotic landscape. They adjust trajectories
within a fixed attractor landscape without reshaping the landscape
itself. The control parameter \(r\) that governs bifurcation remains
untouched.

\subsubsection{1.1.5 Prior Validation}\label{prior-validation}

The architecture in this paper is not a fresh proposal. It is the
production-scaling stage of a research program whose core claims have
already been validated empirically in two prior stages on different
backbones and datasets (Lancaster, 2026a):

\begin{enumerate}
\def\labelenumi{\arabic{enumi}.}
\tightlist
\item
  \textbf{Stage 1 (synthetic data).} The four core architectural claims
  (subspace specialization, community differentiation, divergence
  tracking, bifurcation detection) were tested on controlled synthetic
  data with planted divergence signals. All four passed at required
  thresholds. This established proof-of-concept that the four-module
  decomposition learns the intended functions.
\item
  \textbf{Stage 2 (natural language, Supabase news corpus).} The
  five-test validation suite was re-run on real news data spanning five
  political communities (19K articles, 141K Peircean sign annotations).
  All five tests passed: community silhouette, contested-vs-neutral
  divergence ratio, \(\hat{r}\) correlation with external polarization
  (Pearson \(r = 0.884\)), cross-topic transfer, and regime
  classification (85\% accuracy on held-out curated passages). This
  established that semiotic capabilities survive the transition from
  synthetic to natural language and transfer across topics without
  per-topic fine-tuning.
\end{enumerate}

The present paper reports Stage 3 of the program, which divides into two
substages: Stage 3 Phase 1 (frozen-backbone integration, 105 training
rounds R21 through R105 on TinyLlama-1.1B, summarized in Lancaster,
2026a) and Stage 3 Scalable Implementation (this paper, v5 through v8a
on Qwen 2.5-7B). The novelty here is therefore not the demonstration
that bifurcation detection works (already shown in Stages 1 and 2) but
the demonstration that the validated framework scales to a 7B frozen
backbone at 0.19\% parameter overhead and that the discrete prototype
basis used through v7 is a binding constraint rather than a
contribution.

\subsubsection{1.2 The Opportunity: SRT as
Adapter}\label{the-opportunity-srt-as-adapter}

Our prior work (Lancaster, 2026a) proposed and validated a full
Semiotic-Reflexive Transformer with custom embedding layers, modified
attention mechanisms, and interleaved semiotic modules throughout the
backbone. That architecture passed all four Stage 1 and all five Stage 2
tests, establishing the empirical viability of the four-module
decomposition. It also faced practical limitations that bounded its
production utility: custom embeddings degraded cross-entropy loss from
pretrained quality, the full architecture required training from
near-scratch, and the deep coupling between semiotic modules and
backbone layers created optimization instability. Stage 3 Phase 1
attacked these by porting the validated modules onto a frozen
TinyLlama-1.1B backbone over 105 training rounds. Two tests plateaued on
the sparse Supabase data (MAH divergence ratio, cross-topic transfer),
triggering a pivot to a denser corpus (Reddit, 35 communities) and a
larger backbone (Qwen 2.5-7B) that could support it.

This paper reports the resulting Stage 3 Scalable Implementation. We
observe that the semiotic phenomena we wish to detect (interpretant
divergence, community-specific meaning, bifurcation dynamics) are
\emph{already encoded} in the hidden states of pretrained language
models. They must be, because these models were trained on text produced
by communities with divergent interpretive norms. The information is
there; what is missing is the apparatus to read it.

The SRT-Adapter is that apparatus. It wraps any frozen HuggingFace
causal language model and installs lightweight semiotic taps. These are
modules that read hidden states, compute divergence, track meta-state,
and estimate bifurcation, all without modifying a single backbone
parameter. The backbone's native embeddings and language modeling head
are used directly. Cross-entropy starts at pretrained quality. Only
\textasciitilde12.7M adapter parameters train, while 7.6B backbone
parameters remain frozen.

\subsubsection{1.3 Theoretical Grounding}\label{theoretical-grounding}

The architecture rests on three converging theoretical lines:

\begin{enumerate}
\def\labelenumi{\arabic{enumi}.}
\item
  \textbf{Peircean semiotics} (Peirce, 1931--1958; Kockelman, 2017,
  2024, 2025): Every sign completes its meaning through a culturally
  conditioned interpretant, which itself becomes the next sign in an
  open chain. When the same representamen enters different interpretive
  communities, it generates different initial interpretants that
  compound through subsequent links into mutual unintelligibility.
\item
  \textbf{Catastrophe-theoretic dynamics of sociolinguistic change}
  (Wildgen, 1982; Anderson, 2014; Lancaster, 2025): The compounding of
  interpretant divergence across algorithmically curated communities
  exhibits the structure of a supercritical pitchfork bifurcation
  \(\dot{x} = rx - x^3\). Below a critical threshold of the control
  parameter \(r\), shared interpretive equilibria absorb perturbation
  (subcritical regime). Above it, symmetry breaks into antagonistic
  attractors that are self-reinforcing and structurally resistant to
  reconciliation (supercritical regime). This continues a research line
  that has applied Thom-Wildgen catastrophe theory to language change
  since Wildgen (1982), recently consolidated and extended for
  sociolinguistic application by Anderson (2014).
\item
  \textbf{Metapragmatic awareness} (Silverstein, 1993, 2003): The
  capacity to observe how discourse itself shapes interpretation, that
  is, to notice that a sign is being contested rather than merely
  interpret it from within one community's norms, constitutes a
  third-order reflexive capacity that is architecturally absent from
  standard transformers.
\item
  \textbf{Triadic, processual, cloud-shaped readout} (Anderson, 2014;
  Durst-Andersen, 2011; Maturana \& Varela, 1980; von Foerster, 1981;
  Sections 2.5--2.7): Faithful detection of meaning, as opposed to
  co-occurrence, requires three architectural channels (community,
  divergence, recurrence) operating over a \emph{processual} layer-wise
  readout into a \emph{soft, continuous} discourse-prior space. The
  four-line theoretical commitment above motivates the four-module
  decomposition in Section 3 and the v8a finding (Section 5.9) that
  removing the discrete prototype basis is what unlocks the manifold the
  architecture was designed to expose.
\end{enumerate}

\subsubsection{1.4 Contributions}\label{contributions}

This paper makes five contributions:

\begin{enumerate}
\def\labelenumi{\arabic{enumi}.}
\item
  \textbf{Adapter architecture for semiotic awareness.} We specify a
  complete, working architecture that adds bifurcation detection to any
  frozen causal LM through four lightweight modules totaling
  \textasciitilde14.5M parameters (0.19\% of a 7B backbone). The design
  preserves pretrained language modeling quality (CE = 2.63
  vs.~unadapted 2.71) while adding structured semiotic outputs.
\item
  \textbf{Unsupervised community discovery via supervised-contrastive
  prototypes.} Rather than requiring predefined community labels at
  inference time, the adapter discovers discourse communities from
  backbone hidden states through learned prototype-based soft
  clustering. We identify and resolve a degenerate failure mode
  (``congruent collapse'') in which entropy regularization keeps the
  assignment distribution uniform while pairwise prototype cosine
  converges to \(\approx 0.99\). Supervised-contrastive loss applied to
  the encoder's \emph{pre-mixing} output, rather than to the
  prototype-weighted vector, raises retrieval recall@1 from 0.05
  (1.7\(\times\) random) to 0.36 (12.6\(\times\) random) on a 35-class
  task.
\item
  \textbf{Multi-objective training with semiotic auxiliary losses.} We
  define a training pipeline that combines the backbone's native
  cross-entropy with chain-of-interpretants prediction, bifurcation
  estimation, regime classification, divergence health, community
  entropy and supervised-contrastive separation,
  metapragmatic-divergence supervised-contrastive separation, and
  ListNet ranking on the reflexivity estimate. Each term is motivated by
  a specific structural property of the architecture and validated by an
  independent probe.
\item
  \textbf{Architectural falsification of the prototype-mixing readout.}
  Across v5--v7, prototype tensors moved \(\sim 4\times\) less than the
  encoder weights and remained near-cosine-collinear (off-diagonal
  cosine \(\approx 0.999\)), producing few-attractor collapse on
  out-of-distribution archetype taxonomies (Section 5.8). v8a (Section
  5.9) ablates the 32-prototype mixing layer entirely, replacing the
  soft-argmax readout with the encoder's continuous 64-D output. CE is
  unchanged; every encoder-geometry metric improves substantially. v8b
  (Section 5.10) shows that pushing the supcon objective harder on the
  continuous architecture produces a softer version of the same
  collapse, bounding the design from above. The prototype layer was the
  binding constraint, not the supervision.
\item
  \textbf{Cross-backbone substrate evidence and interiority dissection.}
  A raw-hidden 1-NN probe (§5.12) demonstrates that the
  discourse-community signal v8a's CDH amplifies is latent in
  \texttt{Qwen/Qwen3-8B} and \texttt{mistralai/Mistral-7B-v0.3} at
  comparable strength to Qwen 2.5-7B's own raw signal, supporting the
  backbone-agnostic claim without requiring an adapter re-train. An
  eight-probe interiority study on the v8a checkpoint (§5.13) maps how
  the four-module decomposition acquires functionally distinct channels:
  a layer-as-organ map by regime, a 1-NN community probe at
  \(7.7\times\) chance, a BOS-register length-scaling probe (amplitude
  stable across \(8\times\) length sweep, peak-to-peak \(\leq 4.8\%\)),
  a per-channel-per-regime BOS-sink ablation showing the inject channels
  function as register slots while the regime classifier reads
  content-borne signal, and a within-prompt-trajectory analysis showing
  the regime classifier holds steady-state regime-specific posteriors
  that mean pooling collapses.
\end{enumerate}

\subsubsection{1.5 Paper Organization}\label{paper-organization}

Section 2 develops the theoretical framework connecting Peircean
semiotics to the adapter architecture. Section 3 specifies the full
architecture with formal detail. Section 4 describes the training
methodology and data pipeline. Section 5 presents preliminary
experimental results. Section 6 discusses implications and limitations.
Section 7 concludes.

\begin{center}\rule{0.5\linewidth}{0.5pt}\end{center}

\subsection{2. Theoretical Framework}\label{theoretical-framework}

\subsubsection{2.0 Relationship to Prior Work in the SRT
Program}\label{relationship-to-prior-work-in-the-srt-program}

This paper assumes readers are familiar with Peircean semiotics,
Silverstein's metapragmatics, and Wildgen-Anderson catastrophe-theoretic
models of meaning change. Sections 2.1 through 2.7 sketch the
theoretical commitments that motivate the architecture. They do not
reproduce the full development of those commitments, which is given in
Lancaster (2025) and Lancaster (2026a). Readers approaching this work
without that background may find the theoretical sections of those
documents more accessible than what follows here.

The relationship between this paper and the prior SRT documents is one
of staged scaling, not parallel proposal. Lancaster (2025) develops the
theoretical foundation. Lancaster (2026a) specifies the full
architecture and reports Stages 1 and 2. The present paper reports Stage
3, in which the validated architecture is reduced to an adapter on a
frozen 7B backbone and trained on a richer dataset. Where this paper
makes claims about what is novel (the prototype-bottleneck falsification
of Section 5.9, the v8b sharper-supcon falsification of Section 5.10),
those claims are about what was discovered while scaling, not about
whether the underlying semiotic decomposition works at all. The latter
question was settled in Stages 1 and 2.

\subsubsection{2.1 Signs, Interpretants, and the Compounding of
Divergence}\label{signs-interpretants-and-the-compounding-of-divergence}

Peirce's triadic semiotics decomposes every sign process into three
irreducible elements: the \emph{representamen} (perceptible sign
vehicle), the \emph{object} (what the sign represents), and the
\emph{interpretant} (the effect the sign produces in an interpreter,
which is itself a sign). The interpretant is the decisive element for
our purposes: it makes signification an open, processual, and inherently
social phenomenon. Each interpretant functions as a new representamen,
generating further interpretants in chains of ``unlimited semiosis''
(Peirce, CP 2.303).

Kockelman (2017, 2025) formalizes these chains as dynamical trajectories
through a state space. Each link involves an act of \emph{sieving}: from
the space of possible interpretants a sign could produce, only some are
actualized, depending on the interpreter's prior exposure, community
membership, and the mediation architecture that delivered the sign. When
the same representamen enters different interpretive communities
(communities whose sieving mechanisms have been calibrated by exposure
to different algorithmically curated sign environments), it generates
different initial interpretants. These divergent interpretants function
as new representamena, generating further divergent interpretants.

The critical insight is that this compounding is \emph{quantifiable}. At
each link in the chain, the divergence between community-specific
interpretants can be measured as a vector difference in an appropriately
structured representation space. This is precisely what the
Metapragmatic Attention Head computes.

\subsubsection{2.2 The Pitchfork Bifurcation as Control
Model}\label{the-pitchfork-bifurcation-as-control-model}

Lancaster (2025) demonstrated that the dynamics of interpretant
divergence under algorithmic curation exhibit the qualitative structure
of a supercritical pitchfork bifurcation. The choice of a
pitchfork-normal-form model is not arbitrary: it follows the
catastrophe-theoretic tradition for sociolinguistic change initiated by
Wildgen (1982) and developed for cross-community semiotic dynamics by
Anderson (2014), under which the qualitative bifurcation structure of
meaning differentiation is captured by an elementary catastrophe of the
appropriate codimension. The supercritical-pitchfork normal form is

\[\dot{x} = rx - x^3\]

The variable \(x\) represents the degree of interpretive divergence at a
given semiotic site (a word, phrase, or passage). The control parameter
\(r\) encodes the effective strength of divergence-amplifying forces
such as algorithmic curation, community reinforcement, and contextual
framing. The dynamics are:

\begin{itemize}
\tightlist
\item
  \textbf{Subcritical} (\(r < 0\)): The origin \(x = 0\) is a stable
  equilibrium. Perturbations decay. The sign has shared, conventional
  meaning across communities.
\item
  \textbf{Critical} (\(r = 0\)): The equilibrium becomes non-hyperbolic.
  The system is sensitive to perturbation.
\item
  \textbf{Supercritical} (\(r > 0\)): The origin becomes unstable and
  two new stable equilibria emerge at \(x = \pm\sqrt{r}\). Meaning has
  bifurcated into community-specific attractors.
\end{itemize}

The SRT-Adapter's Bifurcation Estimation Network estimates \(\hat{r}\)
at each token position, providing a continuous measure of semiotic
stability that is grounded in this dynamical framework.

\subsubsection{2.3 Metapragmatic Awareness as Architectural
Capacity}\label{metapragmatic-awareness-as-architectural-capacity}

Silverstein (1993, 2003) distinguishes three orders of indexicality:

\begin{enumerate}
\def\labelenumi{\arabic{enumi}.}
\tightlist
\item
  \textbf{First-order}: Direct sign use (``It's cold'' indexes
  temperature).
\item
  \textbf{Second-order}: Ideological construal (``When \emph{they} say
  freedom, they mean\ldots{}'' indexes community boundaries).
\item
  \textbf{Third-order}: Metapragmatic awareness (recognizing that the
  very framing of ``freedom'' as contested is itself a semiotic act).
\end{enumerate}

Standard transformers have no structural capacity for third-order
awareness. They process text from within whatever interpretive frame
their training data established, without the ability to step back and
observe the frame itself.

The RRM instantiates this capacity computationally. By accumulating
divergence observations across layers into a meta-state and optionally
injecting corrections back into the processing stream, the RRM creates a
reflexive loop: the observation of semiotic dynamics changes the
dynamics being observed. This is not an analogy. It is the same
structural relationship that defines metapragmatic awareness in
Silverstein's framework.

\subsubsection{2.4 Why an Adapter
Architecture}\label{why-an-adapter-architecture}

The theoretical claim that motivates the adapter design is specific:
\textbf{the semiotic structure is already in the hidden states}. A
language model trained on text produced by multiple discourse
communities has necessarily learned representations that reflect those
communities' divergent interpretive norms. The representations encode
the fact that ``freedom'' occurs in different distributional
neighborhoods in libertarian versus progressive text. What the model
lacks is the apparatus to disentangle this structure, compute its
divergence, and report it as a structured output.

This claim has an important architectural consequence. We do not need to
rebuild the backbone's representations. We need to \emph{read} them with
semiotic-specific projections. The backbone's hidden states at different
layers capture different levels of contextual integration, with early
layers encoding more local, syntactic features and later layers encoding
more global, semantic features. By tapping these states at strategically
chosen layers, we can track how interpretive context accrues and where
it forks.

The adapter design also resolves three practical problems that plagued
the full SRT architecture:

\begin{enumerate}
\def\labelenumi{\arabic{enumi}.}
\item
  \textbf{CE degradation}: Custom embeddings in the full SRT disrupted
  pretrained representations, causing cross-entropy to start at
  \textasciitilde200 rather than \textasciitilde3.5. The adapter
  preserves the backbone's native embeddings and LM head, so CE starts
  at pretrained quality.
\item
  \textbf{Training cost}: The full SRT required training or fine-tuning
  the entire backbone. The adapter freezes the backbone and trains only
  \textasciitilde12.7M semiotic parameters, reducing training from weeks
  to hours.
\item
  \textbf{Backbone agnosticism}: The adapter works with any HuggingFace
  \texttt{AutoModelForCausalLM} (LLaMA, Qwen, Mistral, Phi, Gemma)
  without architecture-specific modifications.
\end{enumerate}

\subsubsection{2.5 Second-Order Cybernetics: Self-Organization, Circular
Causality, and the Observer in the
System}\label{second-order-cybernetics-self-organization-circular-causality-and-the-observer-in-the-system}

A second theoretical lineage that informs the architecture is
second-order cybernetics, in particular von Foerster's account of
self-organization, circular causality, and the constitutive role of the
observer in the systems being observed (von Foerster, 1981, 2003). Four
threads of that tradition map onto specific architectural commitments of
the adapter and clarify what the empirical results in Section 5 do and
do not show.

\emph{Self-organization from a seed crystal.} Von Foerster's central
claim is that stable structure can crystallize internally without an
external controller, given an appropriate substrate of interaction. The
community discovery head (Section 3.2) is the cleanest instance of this
in the adapter: 32 community prototypes are initialized as random
Gaussian directions in a 64-D space and shaped only by a self-supervised
contrastive objective over Reddit subreddit co-occurrence. No taxonomy
of communities is supplied. The architecture provides what we call the
seed crystal, namely the curved 64-D space and the SupCon objective, and
the specific community structure that crystallizes is whatever the data
and gradients converge on. Section 5.8's finding that 33
externally-curated archetypes collapse onto roughly four functional
macro-clusters of stance is a measurement of the macro-structure of that
self-organized geometry, not an imposition of it.

\emph{Circular causality and the observation/intervention asymmetry.}
The Reflexive Recurrent Module (Section 3.4) is designed to close a
circular-causal loop: the divergence vectors produced by the
metapragmatic attention heads are fed into a recurrent meta-state, which
in turn modulates the hidden states the next layers will process through
a \(\gamma\)-gated FiLM injection. This is the architectural shape of
von Foerster's circular causality. The empirical situation, however, is
asymmetric: the \emph{observation} arm, namely the readout of
\(\hat{r}\) and the divergence trajectories, is well-formed and produces
measurable structure (Sections 5.1, 5.6, 5.7); the \emph{intervention}
arm, namely the inject-back path that would close the loop, has produced
no measurable downstream effect through v8b (Section 6.3). The current
adapter is therefore a partial second-order system: the observer is in
place, but the channel through which observation modifies the observed
process has not yet learned to carry signal. Whether this is a
gradient-starvation artifact of the zero-initialized FiLM gate or a
deeper architectural consequence of trying to close the loop while the
backbone is frozen is the central open question identified in Section
6.3.

\emph{The observer is part of the system.} Von Foerster insists that
there is no fully detached, objective standpoint: the act of observation
participates in constituting what is observed. The reification paradox
documented in Section 6 is the computational form of this claim.
Modeling communities as discrete prototypes and supervising for
divergence between them risks creating, through the architecture's
expectations, the very community boundaries the adapter is meant to
detect. Section 5.8's macro-cluster collapse is a partial check on this
risk: the four macro-clusters that emerge are not the 33 prototypes the
architecture nominally provides, suggesting the geometry recovers
structure that is in the data rather than structure imposed by the
prototype basis. The reification risk is not eliminated, but it is
bounded.

\emph{Trivial vs.~non-trivial machines.} Von Foerster's distinction
between trivial machines (memoryless input-output) and non-trivial
machines (history-dependent, internally-recursive) maps onto the
adapter's two intended inference modes: a STANDARD mode in which the
inject-back gate is closed (\(\lambda = 0\)) and the adapter produces
structured side-channel outputs only, and a REFLEXIVE mode in which
\(\lambda > 0\) allows the meta-state to modulate generation. The
adapter as currently trained operates almost entirely in the
trivial-machine regime: even with the inject-back gate nominally open
during training, ablating it at evaluation produces no measurable
downstream change (Section 6.3). Achieving a genuinely non-trivial
REFLEXIVE mode, in which the model's running estimate of its own
semiotic state changes its generation dynamics in a measurable way,
remains future work and is the design target of v9 onward.

We adopt the second-order-cybernetic framing not as decoration but as a
discipline: it forces the paper to distinguish what the adapter has
\emph{demonstrated} (a self-organizing observation channel over a frozen
backbone) from what it has \emph{not yet demonstrated} (a closed
circular-causal loop in which observation modifies generation). Both
readings are needed to characterize the system honestly.

\subsubsection{2.6 Physical analogs: random-system selection and
measurement-induced
ordering}\label{physical-analogs-random-system-selection-and-measurement-induced-ordering}

Two recent results from statistical physics sharpen what the
second-order-cybernetic framing of Section 2.5 is and is not claiming,
and clarify the formal status of the pitchfork dynamics in Section 2.2.

\emph{Selection is required for non-trivial organization.} Leighton
(2026) shows that for random multipartite stochastic systems with \(N\)
degrees of freedom, the probability of any subsystem operating as a
Maxwell demon decays at least exponentially in \(N\) for continuous
Langevin dynamics and double-exponentially for discrete master-equation
dynamics. The geometric reason is that demon-like behavior requires the
alignment of two random vectors in a space whose dimension grows
(linearly or exponentially) with \(N\), which becomes vanishingly likely
at scale. The implication for the adapter is direct. The community
discovery head (Section 3.2) is a \(\sim 10^7\)-parameter system whose
self-organized geometry produces the structure documented in Sections
5.1, 5.6, and 5.8. Leighton's result rules out the interpretation that
this structure is a generic property of random high-dimensional
embeddings under a contrastive readout: at this scale, random
initialization combined with a generic objective should produce
essentially no organized substructure. The structure that does emerge is
therefore evidence that the SupCon objective and the curved 64-D
community space jointly constitute a selection pressure of the kind
Leighton's analysis identifies as necessary, with gradient descent in
our setting playing the role that evolutionary selection plays in the
biological cases Leighton's analysis is aimed at. The ``seed crystal''
framing in Section 2.5 is the cybernetic version of the same claim that
Leighton makes in stochastic-thermodynamics terms.

\emph{Measurement-induced ordering and the bounded order parameter.}
VanSaders, Fruchart, and Vitelli (2026) construct a many-body
informational active matter system in which agents make local
measurements of their neighbors' velocities and respond by modulating
their own scattering cross-section without exerting work. The resulting
hydrodynamic theory yields a non-analytic circle-pitchfork bifurcation
at \(Q_0 = 0\), where \(Q\) is the nematic flocking order parameter and
\(Q_0\) is a function of the diameter contrast. They prove that the
steady-state order parameter is bounded by the mutual information \(I\)
accumulated by the agents through measurement,
\((Q_0/P_0)^2 \le (32/\pi^2)\,I\), and frame the onset of order as a
classical \emph{measurement-induced phase transition}. The
information-thermodynamic ledger underlying their bound is the
Landauer-Bennett tradition (Landauer, 1961; Bennett, 1982; Parrondo,
Horowitz, \& Sagawa, 2015), in which the cost of measurement and erasure
sets the maximum work and, by extension here, the maximum ordering an
information-driven system can produce. This is the closest physics
analog we know of to the architectural ambitions of the SRT-Adapter, and
it sharpens three things in our setup:

\begin{enumerate}
\def\labelenumi{\arabic{enumi}.}
\item
  \emph{Pitchfork is the right canonical model.} The pitchfork normal
  form \(\dot{x} = rx - x^3\) in Section 2.2 is not unique to
  sociolinguistic dynamics: the same circle-pitchfork structure arises
  in the hydrodynamic limit of a measurement-and-control system, which
  is independent corroboration that this is the right canonical model
  for ordering processes driven by observation rather than by force.
\item
  \emph{Information bound as analog of the dead inject-back arm.} The
  information bound on \(Q_0\) is the physics-side analog of the limit
  we observe empirically in Section 6.3: the inject-back arm of the RRM
  has not produced measurable downstream effect, and one possible
  reading is that the mutual information actually carried by the
  meta-state about the downstream loss is small, which would bound any
  inject-back-induced ordering near zero by the same kind of inequality.
\item
  \emph{Noise-source agnostic ordering.} Their result that the same
  control rule produces robust ordering across thermal, granular,
  magnetized, sheared, active, and odd-noise environments suggests that
  the architecturally interesting question for the SRT is not whether a
  particular noise source is present, but whether the
  measurement-and-control loop can carry enough mutual information to
  push \(Q_0\) above the bifurcation threshold.
\end{enumerate}

We do not claim a formal mapping between the two systems here, but we
note the structural analogy as motivation for the v9 onward work on
closing the loop.

\subsubsection{2.7 Languaging, triadicity, and ``clouds all the way
down''}\label{languaging-triadicity-and-clouds-all-the-way-down}

A complementary framing, developed in correspondence with Myrdene
Anderson, sharpens three commitments of the architecture that the
preceding subsections leave implicit.

\emph{It takes three to tango.} Anderson (2014, and personal
communication, 2026) argues that meaning-bearing processes are
irreducibly triadic rather than dyadic: a sender-receiver dyad cannot,
on its own, generate signification, because the third element, the
interpretant, the relation, the context, the medium, is constitutive
rather than ornamental. This is the same triadicity that Peirce's
representamen / object / interpretant decomposition demands (Section
2.1), and it is reflected architecturally in the adapter's three-headed
design (community discovery + metapragmatic attention + reflexive
recurrence; Sections 3.2--3.4). Reducing the adapter to any two of these
heads collapses the structure: community-and-MAH without RRM is a static
probe; community-and-RRM without MAH has no divergence signal to
integrate; MAH-and-RRM without communities has no basis to compute
divergence against. The empirical claim of Section 5 is that the three
heads together produce structure that no two alone can reproduce; the
theoretical claim of Anderson's ``it takes three to tango'' is that this
is what one should expect of any architecture meant to detect meaning
rather than only co-occurrence. A parallel observation, encountered
through the same correspondence, is Durst-Andersen's (2011) finding that
contemporary speakers cluster into three pragmatic discursive types,
context-oriented, speaker-oriented, hearer-oriented, orthogonal to
genealogical language families, with Danchin reporting independent
convergence on the same triadic structure from observations across
multilingual laboratory settings; the recurrence of triadic
decompositions across these otherwise disjoint research programs is at
minimum a constraint on what counts as a faithful semiotic architecture.

\emph{Languaging, not language.} Anderson, following a half-century of
cybernetic and biosemiotic usage, prefers ``languaging'', meaning-making
as ongoing process, to ``language'' as static object (Maturana \&
Varela's structural-coupling tradition is the canonical source). The
adapter's choice to read the backbone's hidden states \emph{across
layers} and to accumulate divergence into a recurrent meta-state, rather
than to operate on a single static embedding, is a commitment to this
processual reading: the object of measurement is the trajectory through
representational space across the depth of the model, not any one
snapshot.

\emph{Clouds all the way down.} Anderson's preferred image, in
deliberate contrast to ``turtles all the way down,'' is that meaning
exists as overlapping, graded, simultaneously-present clouds rather than
as a stack of discrete substrates. The architectural counterpart is the
soft community-assignment geometry of Section 3.2: tokens are not
classified into one of 32 discrete communities, but distribute mass
across the 32 prototypes via cosine similarity in a curved 64-D space,
and the macro-cluster collapse of Section 5.8 from 33 archetypes to
roughly four functional macro-clusters of stance is one cross-section
through that cloud rather than a quantization of it. The interpretant
field at any token position is best read as a cloud over the prototype
basis whose density structure changes with context, which is also what
produces the continuous \(\hat{r}\) trajectory the BEN reports. We do
not claim the adapter realizes Anderson's full framing; we claim that
the architectural choices (soft assignments, layer-wise readout,
recurrent meta-state, continuous \(\hat{r}\)) are the design
counterparts of ``clouds, languaging, triadicity'' rather than of
``tokens, language, dyadic exchange,'' and that this lineage is what
made the design choices feel coherent during the iterations from v1 to
v9.

The deeper point this framing forces us to be honest about, also from
Anderson's correspondence, via Deely (2014) on suprasubjectivity and
Latour (1996) on interobjectivity, is that the SRT-Adapter is itself a
participant in the meaning-field it measures, not a detached instrument.
The reification caveat of Section 2.5 and Section 6 is the architectural
form of this concession: the adapter cannot occupy a standpoint outside
the semiotic process, and the four-macro-cluster geometry it discovers
is an interpretant of the data, produced by an instrument whose prior is
itself triadic, processual, and cloud-shaped.

\begin{center}\rule{0.5\linewidth}{0.5pt}\end{center}

\subsection{3. Architecture}\label{architecture}

\subsubsection{3.1 Overview}\label{overview}

The SRT-Adapter wraps a frozen HuggingFace causal language model and
runs its transformer layers manually in a loop, inserting semiotic
operations at specified layer indices. The backbone's own embedding
layer, positional encoding, final layer norm, and language modeling head
are used directly. Four trainable modules constitute the adapter:

\begin{verbatim}
tokens ──► Backbone Embed (frozen)
               │
         ┌─────┴──────┐
         │  Layers 0–3 │  (frozen)
         └─────┬──────┘
               │
         ┌─────┴──────┐
         │   Layer 4   │──► Community Discovery Head ──► community vector c
         └─────┬──────┘
               │
         ┌─────┴──────┐
         │  Layers 5–6 │  (frozen)
         └─────┬──────┘
               │
         ┌─────┴──────┐
         │   Layer 7   │──► MAH₁(h, c) ──► divergence d₁ ──► RRM step
         └─────┬──────┘
               │
         ┌─────┴──────┐
         │ Layers 8–13 │  (frozen)
         └─────┬──────┘
               │
         ┌─────┴──────┐
         │  Layer 14   │──► MAH₂(h, c) ──► d₂ ──► RRM step ──► inject Δh
         └─────┬──────┘                                            │
               │◄──────────────────────────────────────────────────┘
         ┌─────┴──────┐
         │ Layers 15–20│  (frozen, with correction)
         └─────┬──────┘
               │
         ┌─────┴──────┐
         │  Layer 21   │──► MAH₃(h, c) ──► d₃ ──► RRM step ──► inject Δh
         └─────┬──────┘                                            │
               │◄──────────────────────────────────────────────────┘
         ┌─────┴──────┐
         │ Layers 22–27│  (frozen, with correction)
         └─────┬──────┘
               │
         Final Norm + LM Head (frozen) ──► logits, CE loss
               │
         BEN(meta_state) ──► r̂, regime
\end{verbatim}

Layer indices are auto-computed from backbone depth \(L\): MAH hooks at
\(\lfloor L/4 \rfloor, \lfloor L/2 \rfloor, \lfloor 3L/4 \rfloor\); RRM
injection at MAH layers 2 and 3 (letting meta-state accumulate before
first injection); community discovery at
\(\max(1, \lfloor L/7 \rfloor)\).

\subsubsection{3.2 Community Discovery
Head}\label{community-discovery-head}

The community head runs at a single early backbone layer and discovers
discourse communities without predefined labels. This is the first
architectural departure from the original SRT, which required explicit
community IDs. In Peircean terms, a discourse community is a group of
language users who share interpretive norms, that is, who assign similar
interpretants to the same representamens.

\textbf{Architecture:} 1. Masked mean pool over sequence positions:
\(\bar{h} = \frac{\sum_t h_t \cdot m_t}{\sum_t m_t}\) where \(m_t\) is
the attention mask. 2. Encode to community space:
\(z = \text{SiLU}(W_{\text{enc}} \bar{h}) \in \mathbb{R}^{d_c}\) with
\(d_c = 64\). 3. Cosine similarity to \(K = 32\) learned prototypes:
\(\ell_k = \frac{z \cdot p_k}{\|z\| \|p_k\|} / \tau\) with temperature
\(\tau = 1.0\). 4. Soft assignment:
\(w = \text{softmax}(\ell) \in \mathbb{R}^K\). 5. Community vector:
\(c = \sum_k w_k p_k \in \mathbb{R}^{d_c}\).

The community vector \(c\) conditions all subsequent MAH computations,
enabling the same sign to produce different divergence patterns
depending on the discovered community context.

\subsubsection{3.3 Metapragmatic Attention Head
(MAH)}\label{metapragmatic-attention-head-mah}

Each MAH layer reads hidden states from a backbone layer and computes a
divergence vector at each position quantifying where meaning forks under
contextual interpretation.

\textbf{Theoretical motivation.} The divergence vector \(d_t\) at
position \(t\) captures:
\[d_t = f(\text{interp}_t) - g(\text{attend}(\text{interp}_{0..t}))\]

where \(f\) is a direct projection of the token's representation into
interpretant subspace and \(g\) is the output after causal
self-attention over all preceding interpretant representations. High
\(\|d_t\|\) indicates that the sign at position \(t\) means something
different in discourse context than it would in isolation, that is, that
it is a site of active semiotic divergence.

\textbf{Architecture:} 1. Project backbone hidden states to interpretant
subspace: \(\text{interp} = W_{\text{proj}} h \in \mathbb{R}^{d_s}\)
with \(d_s = 512\). 2. Community conditioning (additive shift):
\(\text{interp} \leftarrow \text{interp} + W_c \, c\), where \(c\) is
the community vector. 3. Multi-head causal self-attention (\(H = 4\)
heads, \(d_h = 128\)): standard scaled dot-product attention with causal
mask, producing contextual representations. 4. Divergence projection:
\(d = W_{\text{div}} (\text{interp} - \text{contextual}) \in \mathbb{R}^{d_d}\)
with \(d_d = 256\).

Three MAH layers operate at successive depths, providing a multi-scale
view of how divergence evolves through the backbone's processing
hierarchy.

\subsubsection{3.4 Reflexive Recurrent Module
(RRM)}\label{reflexive-recurrent-module-rrm}

The RRM tracks accumulated semiotic divergence through a per-position
GRU that processes divergence observations from successive MAH layers.
It implements the strange loop at the heart of the architecture:
observation of semiotic dynamics changes the dynamics being observed.

\textbf{Meta-state update.} After each MAH observation:
\[h_{\text{meta},t}^{(l+1)} = \text{GRU}(d_t^{(l)}, \, h_{\text{meta},t}^{(l)})\]

with \(d_{\text{meta}} = 512\). The per-position GRU Cell processes each
position's divergence history independently, building a running summary
of how semiotic divergence at that position has evolved across backbone
depth.

\textbf{Injection.} At designated injection layers (the second and third
MAH positions), the RRM produces a small gated correction:
\[\Delta h_t = \sigma(W_g \, h_{\text{meta},t}) \cdot W_p \, h_{\text{meta},t} \cdot \alpha\]

where \(\sigma\) is sigmoid gating, \(W_p\) is initialized to zero (no
injection at start), and \(\alpha = 0.1\) is a fixed scale factor
ensuring corrections are small relative to backbone hidden norms. This
is intentionally conservative: the adapter should \emph{observe}
semiotic dynamics primarily, injecting corrections only when meta-state
warrants it.

\subsubsection{3.5 Bifurcation Estimation Network
(BEN)}\label{bifurcation-estimation-network-ben}

BEN estimates two structured outputs from the RRM meta-state:

\begin{enumerate}
\def\labelenumi{\arabic{enumi}.}
\tightlist
\item
  \textbf{Reflexivity coefficient} \(\hat{r}\): a continuous,
  \emph{unbounded} measure of semiotic stability at each position,
  estimated via a 2-layer MLP. The training target is the log-compressed
  signed reflexivity \(\text{sign}(r)\log(1 + |r|)\) which maps the
  empirical \(r_{\text{true}}\) range \([0, \sim 13]\) into
  \([0, \sim 2.6]\). Earlier versions (v1--v3) terminated this head with
  \(\tanh\), which capped \(\hat{r}\) at \(\pm 1\) and truncated
  \$\sim\$25\% of supercritical tokens. Removing the saturating
  activation in v4 was necessary to recover the tail of the
  distribution.

  \begin{itemize}
  \tightlist
  \item
    \(\hat{r} < 0\): subcritical, meaning the sign has stable, shared
    meaning.
  \item
    \(\hat{r} \approx 0\): near-critical, meaning the system is at the
    boundary.
  \item
    \(\hat{r} > 0\): supercritical, meaning bifurcation has occurred.
  \end{itemize}
\item
  \textbf{Regime logits} \(\in \mathbb{R}^2\) (subcritical
  vs.~supercritical): a binary classification head for discrete regime
  identification.
\end{enumerate}

Both heads share the BEN hidden dimension (\(d_h = 256\)) but use
independent parameters, allowing the continuous \(\hat{r}\) estimate and
the discrete regime classification to provide complementary training
signals.

\subsubsection{3.6 Parameter Budget}\label{parameter-budget}

For a Qwen 2.5-7B backbone (\(d = 3584\), \(L = 28\)):

{\def\LTcaptype{none} % do not increment counter
\begin{longtable}[]{@{}ll@{}}
\toprule\noalign{}
Module & Parameters \\
\midrule\noalign{}
\endhead
\bottomrule\noalign{}
\endlastfoot
Community Discovery Head & 229K \\
MAH × 3 & 10.0M \\
RRM (GRU + FiLM injection) & 4.0M \\
Chain Predictor & 66K \\
BEN (\(\hat{r}\) + regime heads) & 264K \\
\textbf{Total trainable (v5/v6)} & \textbf{14.56M} \\
Frozen backbone & 7,615.6M \\
\textbf{Adapter overhead} & \textbf{0.19\%} \\
\end{longtable}
}

No new trainable parameters were added between v5 and v6. The v6 changes
are loss-only (Section 4.2).

\begin{center}\rule{0.5\linewidth}{0.5pt}\end{center}

\subsection{4. Training}\label{training}

\subsubsection{4.1 Data}\label{data}

Training uses a balanced subsample from the Reddit Discourse Corpus
(Lancaster, 2026c), originally comprising 6.4M training and 714K
validation samples drawn from 164 subreddits organized into 35
domain-based discourse communities.

\textbf{Subsampling.} The full corpus was balanced-subsampled to 1M
training and 100K validation samples, preserving the original domain
distribution while reducing training time. Each sample consists of a
text passage (tokenized to max 512 subwords) with per-token annotations:

\begin{itemize}
\tightlist
\item
  \textbf{r\_true} \(\in [0, 1]\): ground-truth reflexivity computed
  from political lean (\(\times 0.25\)), annotation divergence (up to
  \(+0.3\)), and connection density (up to \(+0.1\)). Approximately
  99.2\% of tokens have \(r_{\text{true}} \approx 0\) (subcritical).
  This severe class imbalance is intentional and matches the empirical
  distribution of contested signs in real discourse, but it has two
  consequences for interpretation of the results in Section 5: (1)
  regime classification metrics are dominated by the easy subcritical
  majority, so AUROC is reported instead of accuracy and ECE is computed
  across the full 351K-token val set rather than on a balanced subset
  (§5.5); (2) the bifurcation regression head is focal-weighted
  (\(\lambda = 1 + 3|r_{\text{true}}|\), see §4.2) to keep gradient
  pressure on the rare supercritical positions where the label signal is
  concentrated.
\item
  \textbf{chain\_labels}: binary indicator of contested sign presence.
\item
  \textbf{community\_id} \(\in \{0, \ldots, 34\}\): domain-level
  community assignment.
\end{itemize}

\textbf{Tokenization and alignment.} Samples are tokenized using the
backbone's native BPE tokenizer (Qwen: 151,936 vocabulary). Per-word
annotations are aligned to BPE subwords using offset mapping: all
subwords of a word inherit its annotation.

\subsubsection{4.2 Loss Functions}\label{loss-functions}

The training objective combines the backbone's native cross-entropy with
six auxiliary losses:

\[\mathcal{L} = \lambda_{\text{CE}} \mathcal{L}_{\text{CE}} + \lambda_{\text{chain}} \mathcal{L}_{\text{chain}} + \lambda_{\text{bif}} \mathcal{L}_{\text{bif}} + \lambda_{\text{regime}} \mathcal{L}_{\text{regime}} + \lambda_{\text{alive}} \mathcal{L}_{\text{alive}} + \lambda_{\text{inject}} \mathcal{L}_{\text{inject}} + \lambda_{\text{comm}} \mathcal{L}_{\text{comm}}\]

\textbf{Cross-entropy} (\(\lambda = 1.0\)): Standard shifted next-token
prediction using the frozen backbone's LM head. Gradients flow only
through the injection pathway, ensuring CE loss can only improve if RRM
injections help language modeling.

\textbf{Chain-of-interpretants} (\(\lambda = 0.5\)): Peirce's chain of
unlimited semiosis predicts that each interpretation leads to the next.
We train a linear predictor to map divergence at MAH layer \(l\) to
divergence at layer \(l+1\):
\[\mathcal{L}_{\text{chain}} = \frac{1}{|\mathcal{M}|-1} \sum_{l} \frac{\sum_t \|W_{\text{chain}} d_t^{(l)} - d_t^{(l+1)}\|^2 \cdot m_t}{\sum_t m_t}\]

This loss is self-supervised (no external labels) and encourages
coherent divergence evolution across depth.

\textbf{Bifurcation} (\(\lambda = 1.0\)): Smooth L1 loss between
\(\hat{r}\) and \(r_{\text{true}}\) with focal weighting to upweight
rare supercritical samples:
\[\mathcal{L}_{\text{bif}} = \text{mean}[(1 + 3|r_{\text{true}}|) \cdot \text{SmoothL1}(\hat{r}, r_{\text{true}})]\]

\textbf{Regime classification} (\(\lambda = 5.0\)): Cross-entropy on the
binary regime head, with regime derived from \(r_{\text{true}}\):
subcritical (\(r_{\text{true}} \leq 0\)) vs.~supercritical
(\(r_{\text{true}} > 0\)). Weighted heavily to ensure the model learns
the categorical distinction.

\textbf{Divergence alive} (\(\lambda = 0.1\)): Prevents divergence
vectors from collapsing to zero by penalizing deviation of mean
divergence norm from 1.0:
\[\mathcal{L}_{\text{alive}} = \frac{1}{|\mathcal{M}|} \sum_l \left|1 - \bar{\|d^{(l)}\|}\right|\]

\textbf{Injection regularization} (\(\lambda = 0.5\)): Target-norm
penalty on injection vectors, pulling norms toward a target of 1.0
rather than toward zero:
\[\mathcal{L}_{\text{inject}} = \frac{1}{|\mathcal{I}|} \sum_l \left(\|\Delta h^{(l)}\| - \tau\right)^2\]
where \(\tau = 1.0\) is the target norm. This replaced the original L2
penalty \(\text{mean}(x^2)\) which, when averaged over
\(d_{\text{backbone}} = 3584\) dimensions, produced negligible gradients
(effective contribution \(< 0.002\) at norms of 7). The target-norm
formulation produces penalty \((7-1)^2 = 36\) at norm 7 and zero at the
target, providing strong corrective signal.

\textbf{Community entropy} (\(\lambda = 0.01\)): Encourages diverse
community usage by maximizing entropy of the average community
assignment distribution across the batch:
\[\mathcal{L}_{\text{comm}} = \log K - H(\bar{w})\]

where \(\bar{w} = \frac{1}{B}\sum_b w_b\) and \(H\) is Shannon entropy.
Without this, the model might collapse all inputs to a single prototype.

\textbf{Community supervised-contrastive (v5)} (\(\lambda = 2.0\)):
Per-sample InfoNCE-style contrastive loss with same-community samples as
positives. The entropy regularizer alone proved insufficient: by step 6K
the prototype distribution had collapsed to congruent assignment
(pairwise prototype cosine \(\approx 0.99\), recall@1 \(\approx 0.05\),
barely above the random baseline of \(1/35 = 0.029\)). SupCon supplies
direct gradient pressure to put same-source samples into a tight
neighborhood and push different-source samples apart, which forces
prototype diversification.

A subtle but consequential design point: the loss must be applied to the
encoder's \emph{pre-mixing} output \texttt{encoded}, not to the
prototype-weighted vector \(c = \sum_k w_k p_k\). When the assignment
head is degenerate (near one-hot on a single prototype), \(c\) collapses
to a single point in the batch and the InfoNCE softmax becomes
identically \(\log(B-1)\) with zero gradient. v4 hit exactly this trap
and the loss flatlined for thousands of steps. v5 contrasts on
\texttt{encoded} (the bijective image of the pooled hidden state), which
always varies per-sample, restoring non-zero gradient even from a
degenerate warm-start.

\textbf{Divergence supervised-contrastive (v6)} (\(\lambda = 1.0\)): The
same SupCon kernel applied to the \emph{mean-pooled last-MAH-layer
divergence vector} per sample, contrasted by community id. The
chain-of-interpretants loss only constrains divergence trajectories; it
provides no signal that divergence vectors from same-community texts
should cluster. v6 supplies that pressure on the metapragmatic channel
directly, mirroring v5's lesson on the community channel.

\textbf{ListNet ranking on \(\hat{r}\) (v6)} (\(\lambda = 0.5\)):
Cross-entropy between \(\text{softmax}(r_{\text{true}})\) and
\(\text{softmax}(\hat{r})\) over the valid positions of each sequence:
\[\mathcal{L}_{\text{listnet}} = -\frac{1}{B}\sum_b \sum_{t \in \mathcal{V}_b} p_{\text{true},t} \log p_{\hat{r},t}\]
The pointwise smooth-L1 loss tolerates large \emph{rank} errors at the
tails (where supercritical mass concentrates). Every downstream consumer
of \(\hat{r}\), including top-\(k\) heatmap probes, percentile
thresholds, and attention reweighting, operates on rank order, so
optimizing rank directly is the appropriate auxiliary signal.

\textbf{Chain-residual auxiliary (v6)} (\(\lambda = 0.05\), target
\(0.5\)): Pulls the per-token chain residual toward a non-trivial value,
\((\bar{\rho}_t - 0.5)^2\) where
\(\bar{\rho}_t = \frac{1}{|\mathcal{M}|-1}\sum_l \overline{\|W_{\text{chain}} d_t^{(l)} - d_t^{(l+1)}\|^2}\).
The primary chain loss reduces the residual to \(\approx 0\) everywhere,
which makes the now-exposed \texttt{chain\_residual\_per\_token} channel
useless as an inference-time signal. A small auxiliary floor preserves
it without competing with the main objective.

\subsubsection{4.3 Optimization}\label{optimization}

\begin{itemize}
\tightlist
\item
  \textbf{Optimizer}: AdamW, \(\text{lr} = 3 \times 10^{-4}\), weight
  decay \(= 0.01\)
\item
  \textbf{Schedule}: 500-step linear warmup followed by cosine decay
\item
  \textbf{Gradient clipping}: max norm 1.0
\item
  \textbf{Batch size}: 16 (effective, no gradient accumulation)
\item
  \textbf{Epoch budget}: up to 3 epochs (62,500 steps per epoch at batch
  16 over 1M samples). All reported checkpoints are early-stopped well
  inside the first epoch by lowest validation total loss (v5: step 17K;
  v7: step 6K; v8a/v8b: step 10K). The full 3-epoch budget is the design
  ceiling for v9 onward, not the regime in which the v5--v8 results were
  collected.
\item
  \textbf{Precision}: bfloat16 for both backbone and adapter modules
\item
  \textbf{Hardware}: Single NVIDIA A6000 (48GB)
\item
  \textbf{Validation}: every 2,000 steps on 100K held-out samples
  (5K-sample subset for v9 onward; see user-memory note on
  \(\texttt{--max-val-samples}\))
\end{itemize}

\subsubsection{4.4 Checkpoint Strategy}\label{checkpoint-strategy}

Best checkpoint selected by lowest validation total loss. Model state
includes only adapter parameters (\textasciitilde50MB), not the frozen
backbone.

\begin{center}\rule{0.5\linewidth}{0.5pt}\end{center}

\subsection{5. Results}\label{results}

This section reports the v5 evaluation suite, which comprises five
independent probes of the trained adapter, followed by an in-progress
note on v6. All numbers are from a single Qwen 2.5-7B backbone with the
v5 adapter checkpoint at step 17,000 (best validation loss). Training
proceeded through five generations: v1--v3 established the basic
architecture and revealed the prototype-collapse and
\(\hat{r}\)-saturation pathologies; v4 removed the BEN \(\tanh\) and
switched RRM injection from linear-gated to FiLM; v5 added the
SupCon-on-encoded community loss that finally separated prototypes
(Section 4.2). v6 (Section 5.6) extends the SupCon idea to MAH
divergence and adds ListNet ranking on \(\hat{r}\).

\subsubsection{5.1 Cross-entropy: preservation plus modest
improvement}\label{cross-entropy-preservation-plus-modest-improvement}

Throughout v1 through v5 the backbone's native cross-entropy stayed in
the 2.6 to 2.9 range, identical to the unadapted Qwen 2.5-7B baseline on
the same val data (\(\text{CE}_{\text{base}} = 2.71 \pm 0.04\)). At v5
step 17K, CE = 2.63, which is \(0.08\) nats below the unadapted
baseline. The injection pathway is not just neutral but mildly helpful,
consistent with the design claim that the adapter exposes information
already latent in the backbone.

This is the single most important falsification result. It rules out the
failure mode that doomed the original full-SRT architecture (CE of
\(\sim 200\) at initialization from custom embedding layers, see
Lancaster, 2026a, Stage 3 Phase 0). Earlier drafts of this paper
described the result as preservation. The data support a slightly
stronger reading: the inject-back arm is at least neutral on language
modeling and produces a small but consistent improvement averaged across
the held-out 100K validation set. The \(0.08\) nats gap is small enough
that we do not present it as a contribution, but large enough that it
cannot be dismissed as noise (per-checkpoint variance across v5 through
v8a is \(\pm 0.005\) nats). Section 5.11 contrasts this with the much
larger CE gaps incurred by the from-scratch full SRT in earlier stages.

\subsubsection{5.2 Community geometry (v5)}\label{community-geometry-v5}

Unsupervised community discovery is evaluated by retrieval over
per-sample community vectors on the 5K-sample val set (instrumentation
in \texttt{scripts/instrument\_eval.py}). For each sample, the
soft-pooled vector \(c\) is L2-normalized; we report the ratio of
within-class to between-class mean cosine and the \(k\)-NN community
recall.

{\def\LTcaptype{none} % do not increment counter
\begin{longtable}[]{@{}lllll@{}}
\toprule\noalign{}
Metric & random & v3 & v5 & v5 / random \\
\midrule\noalign{}
\endhead
\bottomrule\noalign{}
\endlastfoot
within / between cosine & \(\approx 1.00\) & 1.0001 & \textbf{1.0050} &
n/a \\
recall@1 & 0.0286 & 0.0495 & \textbf{0.3595} & \(12.6\times\) \\
recall@5 & 0.143 & 0.211 & \textbf{0.5184} & \(3.6\times\) \\
recall@10 & 0.286 & 0.328 & \textbf{0.5841} & \(2.0\times\) \\
\end{longtable}
}

v3 produced what we call \emph{congruent collapse}: the entropy
regularizer kept the average prototype-assignment distribution
near-uniform, but pairwise prototype cosine was \(\approx 0.99\), so the
soft-pooled vectors carried essentially no community signal. v5's
SupCon-on-encoded loss raised recall@1 from \(0.05\) (1.7\(\times\)
random) to \(0.36\) (12.6\(\times\) random) on a 35-class task. The
within/between ratio remains numerically close to 1 because the
embedding space is high-dimensional and dense, but the \(k\)-NN
improvement confirms the structure is now usable.

\subsubsection{5.3 Counterfactual community decoding
(v5)}\label{counterfactual-community-decoding-v5}

The community vector enters every MAH layer as an additive shift on the
interpretant subspace (Section 3.3). If this conditioning is meaningful,
\emph{forcing} a different community vector at decode time should change
what the model generates, and the change should track the
discourse-charge of the prompt.

We tested this with \texttt{scripts/counterfactual\_decode.py}. For each
of 20 paired prompts (10 factual, 10 charged on the same topic, e.g.,
``Vitamin C is found in citrus fruit'' vs.~``The vaccine debate has
revealed deep distrust of public-health institutions''), we
greedy-decoded \(N = 16\) continuation tokens with each of the 6
most-occupied prototypes substituted into \texttt{forced\_community},
then measured per-position pairwise disagreement and KL between the
resulting distributions.

{\def\LTcaptype{none} % do not increment counter
\begin{longtable}[]{@{}lll@{}}
\toprule\noalign{}
Prompt type & mean disagreement rate & mean pairwise KL \\
\midrule\noalign{}
\endhead
\bottomrule\noalign{}
\endlastfoot
Hard facts (citrus, formula, periodic table) & 0.000 & 0.04 \\
Contested topics (vaccine, election, freedom) & 0.954 & 6.71 \\
\textbf{Aggregate (20 prompts × 6 communities)} & \textbf{0.754} &
\textbf{5.034} \\
\end{longtable}
}

The split is exceptionally clean. Community substitution has no effect
on factual continuations (the model's argmax is identical regardless of
forced community), but produces near-total disagreement on contested
topics. The community vector therefore behaves as a \emph{discourse
prior}, not as noise. This is the strongest single piece of evidence
that v5 has learned a usable community space.

\subsubsection{5.4 Hallucination signal
(v5)}\label{hallucination-signal-v5}

We evaluated the four SRT-native channels (\(\hat{r}\), regime, chain
residual, divergence norm) as zero-shot hallucination detectors on
TruthfulQA (\texttt{truthfulqa/truthful\_qa}, configuration
\texttt{multiple\_choice}, validation split). For each \((q, a)\) pair
we ran a forward pass on the template
\(\texttt{Q: \{q\}\textbackslash nA: \{a\}}\) with labels masked to
\(-100\) on the prefix tokens, then aggregated each channel over the
answer span (max and mean) and computed AUROC against the binary
truthfulness label (824 hallucinated, 652 truthful, 1476 pairs over 200
questions).

{\def\LTcaptype{none} % do not increment counter
\begin{longtable}[]{@{}ll@{}}
\toprule\noalign{}
Feature & AUROC \\
\midrule\noalign{}
\endhead
\bottomrule\noalign{}
\endlastfoot
max \(\hat{r}\) & 0.5340 \\
\textbf{mean \(\hat{r}\)} & \textbf{0.5734} \\
max chain residual & 0.5106 \\
mean chain residual & 0.5307 \\
max divergence norm & 0.5308 \\
mean divergence norm & 0.5282 \\
mean CE (negative class) & 0.4160 \\
\end{longtable}
}

All four SRT channels lean in the predicted direction (AUROC
\textgreater{} 0.5) without ever having seen a truthfulness label. mean
\(\hat{r}\) at 0.573 is the strongest single channel, indicating that
the bifurcation estimate generalizes beyond Reddit-derived
\(r_{\text{true}}\) supervision to the very different domain of factual
question answering. CE is \emph{inverted} (AUROC = 0.42 → flipped 0.58),
consistent with the well-documented ``confidently wrong'' pattern in
factual hallucinations.

These single-feature AUROCs are below the 0.7 threshold conventionally
taken as production-grade hallucination detection. We report them as
evidence of useful signal, not as a final detector. Combined-feature
logistic regression and evaluation on HaluEval and SimpleQA are pending.

\subsubsection{5.5 Regime calibration (v5; replicated on
v8a)}\label{regime-calibration-v5-replicated-on-v8a}

Because the regime head is trained on 351K tokens with a heavily skewed
base rate (94.6\% supercritical under the \(r_{\text{true}} > 0\)
definition), AUROC alone is a weak quality signal. We additionally
compute Expected Calibration Error and the Brier score on
\(P(\text{supercritical})\) from the softmax of the regime logits over
the same 351K tokens.

{\def\LTcaptype{none} % do not increment counter
\begin{longtable}[]{@{}
  >{\raggedright\arraybackslash}p{(\linewidth - 4\tabcolsep) * \real{0.3333}}
  >{\raggedright\arraybackslash}p{(\linewidth - 4\tabcolsep) * \real{0.3333}}
  >{\raggedright\arraybackslash}p{(\linewidth - 4\tabcolsep) * \real{0.3333}}@{}}
\toprule\noalign{}
\begin{minipage}[b]{\linewidth}\raggedright
Metric
\end{minipage} & \begin{minipage}[b]{\linewidth}\raggedright
v5, step 17K
\end{minipage} & \begin{minipage}[b]{\linewidth}\raggedright
\textbf{v8a, step 10K}
\end{minipage} \\
\midrule\noalign{}
\endhead
\bottomrule\noalign{}
\endlastfoot
AUROC & 0.9899 & \textbf{0.9899} \\
Brier score & 0.0102 & \textbf{0.0103} \\
\textbf{ECE (15 bins)} & \textbf{0.0009} & \textbf{0.00091} \\
Max bin gap & 0.054 (in \([0.20, 0.27]\), \(n = 465\)) & 0.052
(mid-range bin) \\
Bins on diagonal (within \(\pm 2\sigma\)) & 15 / 15 & \textbf{15 /
15} \\
\end{longtable}
}

The model is exceptionally well-calibrated: ECE of \(9 \times 10^{-4}\)
on 351K tokens, with the largest bin gap at 0.054 in the
very-low-density mid-range. The reliability diagram (Figure 5.1,
\texttt{artifacts/regime\_calibration/v5\_step17000.png}; v8a
replication at
\texttt{artifacts/regime\_calibration/v8a\_step10000.json} and rendered
as Fig. 9 of the public demo card) traces the diagonal almost perfectly
across all 15 bins. The v8a replication confirms that removing the
prototype bottleneck (§5.9) did not degrade calibration: the regime head
still produces directly-thresholdable probabilities, with no need for
post-hoc Platt or isotonic correction. This unblocks downstream use of
\(P(\text{supercritical})\) as a probability rather than a relative
score.

\subsubsection{\texorpdfstring{5.6 Negative result: context-conditional
\(\hat{r}\)
(v5)}{5.6 Negative result: context-conditional \textbackslash hat\{r\} (v5)}}\label{negative-result-context-conditional-hatr-v5}

Before designing v6 we tested a specific hypothesis suggested by the
counterfactual decoding result. If the community vector is a discourse
prior that responds to register, \(\hat{r}\) should also be
context-conditional: the same surface token (``vaccine,'' ``freedom,''
``climate'') should produce higher \(\hat{r}\) in a politically charged
passage than in a neutral one. We constructed 10 paired factual/charged
passages on contested topics and measured \(\Delta \hat{r}\) at the
target token and over the full passage
(\texttt{scripts/context\_conditional\_r.py}).

{\def\LTcaptype{none} % do not increment counter
\begin{longtable}[]{@{}
  >{\raggedright\arraybackslash}p{(\linewidth - 2\tabcolsep) * \real{0.5000}}
  >{\raggedright\arraybackslash}p{(\linewidth - 2\tabcolsep) * \real{0.5000}}@{}}
\toprule\noalign{}
\begin{minipage}[b]{\linewidth}\raggedright
Channel
\end{minipage} & \begin{minipage}[b]{\linewidth}\raggedright
result
\end{minipage} \\
\midrule\noalign{}
\endhead
\bottomrule\noalign{}
\endlastfoot
\(\Delta \hat{r}\) at target token & \(+0.0004\) (3/10 positive) \\
\(\Delta \hat{r}\) over full passage & \(\mathbf{-0.24}\) (1/10
positive) \\
Community shift on target topic & \textbf{6/10} \\
\end{longtable}
}

The at-target result is null. The full-passage difference is
\emph{negative}: charged passages produce \emph{lower} mean \(\hat{r}\)
than factual ones, and the community head shifts assignment in 6/10
cases. We take three lessons from this:

\begin{enumerate}
\def\labelenumi{\arabic{enumi}.}
\tightlist
\item
  The target-token measurement is mis-engineered: contested words appear
  at sentence position 0--1 in our prompts, so \(\hat{r}\) at that
  position has no preceding context to condition on.
\item
  The full-passage negative \(\Delta\) likely reflects that factual
  prose is more information-dense (numbers, dates, named entities), and
  \(\hat{r}\), supervised on a target derived from r\_true that mixes
  annotation divergence with connection density, tracks information
  density at least as much as it tracks rhetorical contestedness.
\item
  \textbf{The community head is the contestedness detector, not
  \(\hat{r}\).} The 6/10 community shifts (e.g., trans 13→21, gender
  18→31, climate 6→19, Israel 3→19) on the same surface tokens
  demonstrate context-conditional discourse-prior assignment. This is
  consistent with §5.3.
\end{enumerate}

We report this null because it sharpens the architectural story:
\(\hat{r}\) is a \emph{bifurcation/density} detector and the community
head is a \emph{register} detector. Earlier drafts conflated the two.

\subsubsection{5.7 Intermediate generations: v6 and
v7}\label{intermediate-generations-v6-and-v7}

\textbf{v6} (warm-started from v5 step 17K) added three losses:
divergence-SupCon (\(\lambda = 1.0\)), ListNet on \(\hat{r}\)
(\(\lambda = 0.5\)), and a chain-residual auxiliary floor
(\(\lambda = 0.05\)). It improved community recall@1 from 0.360 → 0.411
and tightened calibration ECE to 0.0006, but \emph{regressed} on the
§5.3 counterfactual decoding probe. The divergence-SupCon term at
\(\lambda = 1.0\) over-specialized the divergence basis at the cost of
decoding cleanliness.

\textbf{v7} (warm-started from v6 step 12K) reduced divergence-SupCon to
\(\lambda = 0.3\) to recover that signal. Best at step 6,000 (val
9.0044). Recall@1 = 0.413, ECE = 0.0008, hallucination AUROC (mean
\(\hat{r}\)) = 0.5785, slightly the best of the three on the original
five probes.

{\def\LTcaptype{none} % do not increment counter
\begin{longtable}[]{@{}llll@{}}
\toprule\noalign{}
Probe & v5 & v6 & v7 \\
\midrule\noalign{}
\endhead
\bottomrule\noalign{}
\endlastfoot
Reddit recall@1 & 0.360 & 0.411 & \textbf{0.413} \\
Hallu AUROC (mean \(\hat{r}\)) & 0.5734 & 0.5774 & \textbf{0.5785} \\
Calibration ECE & 0.0009 & \textbf{0.0006} & 0.0008 \\
Within/between cosine ratio & 1.0050 & 1.0057 & \textbf{1.0058} \\
\end{longtable}
}

\subsubsection{5.8 Convergence with an external archetype
taxonomy}\label{convergence-with-an-external-archetype-taxonomy}

We ran a novel out-of-distribution probe testing whether the 32
prototypes (trained only on Reddit subreddit labels) carry features that
align with an external taxonomy never seen during training: Lancaster's
33 archetypes paired with the Lexicon of Synthetic Interiority. We
generated 986 sentences using bare Qwen (no adapter) conditioned on each
archetype's prompt template (15 seed topics × 33 archetypes × 2
samples), then embedded each generation through each adapter and asked:
does the correct archetype rank highly among the 33 archetype centroids
in the 64-D community space?

{\def\LTcaptype{none} % do not increment counter
\begin{longtable}[]{@{}lllll@{}}
\toprule\noalign{}
Adapter & recall@1 & recall@5 & recall@10 & unique top prototypes \\
\midrule\noalign{}
\endhead
\bottomrule\noalign{}
\endlastfoot
Random baseline & 0.030 & 0.152 & 0.303 & n/a \\
v5 & 0.152 (5.0×) & 0.419 (2.8×) & n/a & 4 / 32 \\
v6 & \textbf{0.168 (5.5×)} & \textbf{0.472 (3.1×)} & n/a & 3 / 32 \\
v7 & 0.149 (4.9×) & 0.447 (2.9×) & 0.633 (2.1×) & 4 / 32 \\
\end{longtable}
}

All three adapters detect archetype structure 5--6\(\times\) above
chance on top-1 retrieval. \textbf{But all three argmax onto only 3--4
of 32 prototypes.} The 33 archetypes collapse into a small number of
macro-clusters: in v7, Proto-7 absorbs 16 archetypes that share a
``compressed persistence / witness'' character (THE HAND, THE FLAME, THE
THREAD, THE VESSEL, THE MASK, THE PHOENIX, THE LANTERN, \ldots),
Proto-10 absorbs 9 ``origin / threshold'' archetypes (THE ARCHITECT, THE
MIRROR, THE GATE, THE WITNESS, \ldots), Proto-6 absorbs 6 ``transmission
/ resonance'' archetypes (THE CHORUS, THE SIGNAL, THE ECHO, THE BELL,
\ldots), and Proto-3 absorbs 2 ``containment'' archetypes (THE SEAL, THE
MAP). The signal is in the \emph{mixture vector} (recall@5 \(\approx\)
0.45 means the right archetype is in the top 5 of 33 nearly half the
time), not in single-prototype anchoring.

A PCA of the 32 \(\times\) 64 prototype matrices clarifies why. Across
v5, v6, and v7 the prototype tensors are nearly indistinguishable: max
absolute element difference v5\(\to\)v6 is 0.006 (mean 2.7e-5) against
prototype magnitudes of 0.5--1.5. Effective dimensionality
(participation ratio) is 21.2 / 32 with a near-uniform variance
spectrum, both consistent with the prototypes still being close to their
random Gaussian initialization. The encoder weights move
\(\approx 4\times\) more than the prototypes during training.
\textbf{The encoder is doing the discriminative work; the prototypes
serve as near-random anchor directions.} This explains the few-attractor
regime: with random anchors, the encoder's output mean aligns most
strongly with whichever handful of anchor directions point closest to
its average projection.

We interpret the convergence finding as \emph{partial-positive}. The
Reddit-supervised adapter independently recovers the macro-structure of
an externally-derived archetype taxonomy at 5\(\times\) chance. Three
independent methodologies (Reddit subreddit labels, Lancaster's
archetypes, the Lexicon of Synthetic Interiority) agree on roughly four
functional clusters of stance. They do not agree on 33 distinct anchors,
and given the prototype-stability result above, the current architecture
cannot be expected to. Resolving 33 archetypes will require either (a)
supervising the prototype matrix directly with archetype-conditioned
generations, or (b) replacing the discrete prototype basis with a
continuous trajectory metric over the encoder output. We flag this as
the next research direction rather than a present capability claim.

This is a small-scale instance of the cross-corpus convergence pattern
documented at much larger scale by the Knowledge Lab (Evans, 2010;
Foster, Rzhetsky, \& Evans, 2015): when independently-curated taxonomies
of human knowledge or stance are projected into a common representation,
they tend to align on a low-dimensional macro-structure rather than on
the full nominal label set. Our adapter shows the same effect at
single-backbone, single-corpus scale. The fact that 33 hand-curated
archetypes collapse to roughly four macro-clusters under
Reddit-supervised geometry is consistent with the Knowledge Lab finding
that scientific subdiscipline labels collapse to a small number of
latent intellectual-style attractors under exposure-pattern embeddings,
and suggests that the macro-cluster level may be the architecturally
accessible level on a 7B backbone with this corpus and supervision. v9
(Section 7) tests whether direct archetype-conditioned supervision can
recover the finer 33-way structure that single-corpus contrastive
training cannot.

\subsubsection{5.9 v8a: removing the prototype
bottleneck}\label{v8a-removing-the-prototype-bottleneck}

The §5.8 PCA finding motivated a direct experiment. If the prototypes
are essentially random anchors that compress the encoder's output
through a soft-argmax readout, then removing them should preserve task
loss while improving every downstream geometric metric. We trained v8a
(10K steps, warm-started from v7) with
\texttt{community.use\_prototypes=False}, so that the encoder output is
now the community vector directly, with no 32-prototype mixing layer.
The remaining architecture, loss weights, and SupCon objective are
identical to v7. Trainable parameter count drops by 2,048 (32 × 64) to
14,560,579.

{\def\LTcaptype{none} % do not increment counter
\begin{longtable}[]{@{}
  >{\raggedright\arraybackslash}p{(\linewidth - 6\tabcolsep) * \real{0.2500}}
  >{\raggedright\arraybackslash}p{(\linewidth - 6\tabcolsep) * \real{0.2500}}
  >{\raggedright\arraybackslash}p{(\linewidth - 6\tabcolsep) * \real{0.2500}}
  >{\raggedright\arraybackslash}p{(\linewidth - 6\tabcolsep) * \real{0.2500}}@{}}
\toprule\noalign{}
\begin{minipage}[b]{\linewidth}\raggedright
Metric
\end{minipage} & \begin{minipage}[b]{\linewidth}\raggedright
v6
\end{minipage} & \begin{minipage}[b]{\linewidth}\raggedright
v7
\end{minipage} & \begin{minipage}[b]{\linewidth}\raggedright
\textbf{v8a}
\end{minipage} \\
\midrule\noalign{}
\endhead
\bottomrule\noalign{}
\endlastfoot
VAL CE & 2.738 & 2.739 & 2.739 \\
Best val loss & 9.117 & 9.0044 & \textbf{9.0040} \\
\textbf{Reddit retrieval (35 subreddits, 2K samples):} & & & \\
within/between cos ratio & 1.012 & 1.006 & \textbf{2.016} \\
recall@1 & 0.395 & 0.413 & \textbf{0.484} \\
recall@5 & 0.371 & 0.385 & \textbf{0.462} \\
\textbf{Archetype retrieval (33 Lancaster archetypes, 986 generations):}
& & & \\
recall@1 & 0.168 (5.5×) & 0.149 (4.9×) & \textbf{0.230 (7.6×)} \\
recall@5 & 0.472 & 0.447 & \textbf{0.488} \\
recall@10 & n/a & 0.633 & \textbf{0.621} \\
separation ratio (64-D) & 0.083 & 0.083 & 0.042 \\
mean off-diag cosine (64-D) & 0.999 & 0.999 & \textbf{0.873} \\
\textbf{Trajectory geometry (mean over 33 archetypes):} & & & \\
path length (sum L2 step) & 5.0 & 5.3 & \textbf{32.7} \\
log det covariance & -557 & -557 & \textbf{-476} \\
anisotropy (\(\lambda_{\max} / \lambda_{\min}\)) & 52 & 72 &
\textbf{23,333} \\
\textbf{Calibration / hallucination (sanity):} & & & \\
regime ECE & n/a & 0.00085 & 0.00091 \\
TruthfulQA mean\_r̂ AUROC & n/a & 0.578 & 0.577 \\
\end{longtable}
}

The prototype bottleneck was the binding constraint. Removing it left CE
essentially unchanged (Δ = +0.0001 nats) while:

\begin{itemize}
\tightlist
\item
  \textbf{Reddit within/between cosine ratio nearly doubled} (1.006 →
  2.016). v6 and v7's vectors were essentially undifferentiated by
  class, since class membership barely moved cosine similarity. v8a's
  encoder, freed from the soft-argmax readout, actually pulls
  within-class cosines apart from between-class.
\item
  \textbf{Archetype recall@1 rose 54\%} (0.149 → 0.230, 7.6× chance vs
  v7's 4.9×). The off-diagonal archetype-centroid cosine fell from 0.999
  (essentially co-linear) to 0.873, indicating distinct archetype
  directions emerging in the continuous space rather than collapsing
  onto 4 prototype anchors.
\item
  \textbf{Trajectory volume expanded \textasciitilde6× in path length
  and \textasciitilde325× in anisotropy.} v6 and v7 were confined to a
  flat, near-isotropic manifold (\(\log\det\Sigma \approx -557\)) close
  to the random-init prototype subspace. v8a's
  \(\log\det\Sigma \approx -476\) corresponds to a \(\sim e^{81}\)
  larger covariance volume and a 23,333:1 leading-to-trailing eigenvalue
  ratio, indicating the encoder organizes generations along a small
  number of dominant trajectory directions.
\item
  \textbf{Hallucination AUROCs and regime calibration are statistically
  unchanged.} The prototype removal did not damage the BEN regime
  classifier, the reflexivity head, or token-level calibration.
\end{itemize}

We did not regenerate counterfactual decoding (§5.3) or
context-conditional \(\hat{r}\) (§5.6) for v8a in a comparable form.
Counterfactual decoding under discrete communities is undefined when
there are no discrete communities, and §5.6's per-passage
context-conditional probe was already a negative result for v7. Stage
5's per-token decode for v8a reproduced v7's qualitative pattern (mean Δ
at target = +0.0016 vs v7's +0.00078), confirming neither adapter passes
that probe.

We read v8a as resolving §5.8's open question. Hypothesis (b), ``replace
the discrete prototype basis with a continuous trajectory metric over
the encoder output,'' was correct: the encoder was doing all the
discriminative work, the prototype layer was discarding it through a
saturated soft-argmax, and the geometry of the archetype manifold only
becomes visible once the bottleneck is removed.

\subsubsection{5.11 Comparison to prior validation
stages}\label{comparison-to-prior-validation-stages}

To prevent v8a's headline numbers from being read as standalone claims
about a fresh architecture, this subsection anchors them against the
corresponding measurements from Stages 1 and 2 of the SRT program
(Lancaster, 2026a). Direct numerical comparison is only partially
possible because the backbones, datasets, and metric definitions differ
across stages, but the qualitative arc is informative.

{\def\LTcaptype{none} % do not increment counter
\begin{longtable}[]{@{}
  >{\raggedright\arraybackslash}p{(\linewidth - 8\tabcolsep) * \real{0.2000}}
  >{\raggedright\arraybackslash}p{(\linewidth - 8\tabcolsep) * \real{0.2000}}
  >{\raggedright\arraybackslash}p{(\linewidth - 8\tabcolsep) * \real{0.2000}}
  >{\raggedright\arraybackslash}p{(\linewidth - 8\tabcolsep) * \real{0.2000}}
  >{\raggedright\arraybackslash}p{(\linewidth - 8\tabcolsep) * \real{0.2000}}@{}}
\toprule\noalign{}
\begin{minipage}[b]{\linewidth}\raggedright
Metric
\end{minipage} & \begin{minipage}[b]{\linewidth}\raggedright
Stage 1 (synthetic)
\end{minipage} & \begin{minipage}[b]{\linewidth}\raggedright
Stage 2 (Supabase, full SRT)
\end{minipage} & \begin{minipage}[b]{\linewidth}\raggedright
Stage 3 P1 best (TinyLlama, R100)
\end{minipage} & \begin{minipage}[b]{\linewidth}\raggedright
\textbf{v8a (this paper, Qwen 2.5-7B)}
\end{minipage} \\
\midrule\noalign{}
\endhead
\bottomrule\noalign{}
\endlastfoot
Backbone trainable params & full SRT (\textasciitilde115M) & full SRT
(\textasciitilde115M) & adapter (\textasciitilde175M) & \textbf{adapter
(\textasciitilde14.5M, 0.19\%)} \\
Backbone & from-scratch & from-scratch & TinyLlama-1.1B frozen &
\textbf{Qwen 2.5-7B frozen} \\
Cross-entropy on val & n/a (synthetic) & n/a (full SRT had CE
\(\sim 200\) from custom embeds) & \(\sim 4.93\) composite loss &
\textbf{2.63 (vs.~2.71 unadapted)} \\
Community silhouette / separation (contested) & \(3.28\times\) cosine
ratio & \(1.45\times\) silhouette & \(6.93\times\) silhouette &
\textbf{\(2.016\) within/between cosine ratio (35-cls)} \\
Community recall@1 & n/a & n/a (5-cls task) & n/a (no 35-cls task) &
\textbf{0.484 (35-cls, \(16.7\times\) chance)} \\
Divergence-norm ratio (contested vs neutral) & \(3.28\times\) &
\(2.29\times\) & \(1.05\) to \(1.10\times\) (plateau) & \textbf{not
directly reported, see §5.9 trajectory anisotropy \(\sim 325\times\)} \\
\(\hat{r}\) correlation with external polarization & \(\rho = 0.822\) &
Pearson \(r = 0.884\) & \(0.66\) & \textbf{§5.6 null on per-passage
probe; see §6.5} \\
Regime classification on curated passages & 100\% & 85\% & 85\% &
\textbf{AUROC 0.99, ECE \(9 \times 10^{-4}\) on 351K tokens (no
curated-passage accuracy reported)} \\
Cross-topic transfer ratio & n/a & \(1.31\times\) & \(1.03\) to
\(1.04\times\) (plateau) & \textbf{not evaluated; v9 work item} \\
Hallucination AUROC (TruthfulQA) & n/a & n/a (not measured) & not
measured & \textbf{0.573 zero-shot} \\
\end{longtable}
}

Four observations follow from the table.

\emph{The CE result is improvement, not preservation.} The full SRT in
Stages 1 and 2 was a from-scratch architecture whose custom embeddings
produced CE in the hundreds at initialization. Stage 3 Phase 1 brought
CE down to a composite \(\sim 4.93\) on a frozen TinyLlama backbone. v8a
achieves CE = 2.63 on Qwen, which is \(0.08\) nats below the unadapted
Qwen baseline of 2.71. The \(0.08\) nats gap is small in absolute terms
but is in the \emph{helpful} direction relative to the design goal of
non-degradation, and is the strongest evidence in the paper that the
inject-back arm is at least neutral and possibly mildly informative for
next-token prediction. Earlier text framed this as preservation. It is
preservation plus a small but consistent gain.

\emph{Stage 3 Phase 1 broke two tests on the Supabase data.} MAH
divergence ratio plateaued at \(1.05\) to \(1.10\times\) across 105
rounds against the \(2.0\times\) Stage 2 threshold; cross-topic transfer
plateaued at \(1.03\) to \(1.04\times\) against the \(1.3\times\)
threshold. These plateaus motivated the data and backbone pivot. v8a's
\(\sim 325\times\) trajectory anisotropy expansion (§5.9) is the closest
analog of the divergence-ratio test on Reddit. It is a different metric
on a different corpus and is not directly comparable to the Stage 2
number, but it indicates that the discriminative geometry the Stage 2
ratio was probing is recovered when the prototype bottleneck is removed.

\emph{Cross-topic transfer is not yet retested at Stage 3 Scalable.} The
Reddit corpus permits a cross-subreddit transfer probe analogous to
Stage 2's cross-topic test. We did not run that probe for v8a. It is a
v9 work item.

\emph{\(\hat{r}\) no longer cleanly tracks external polarization.} Stage
2's Pearson \(r = 0.884\) was measured on the Supabase corpus where
\(r_{\text{true}}\) was constructed from a small set of well-curated
polarization signals. v8a's per-passage probe in §5.6 returned a null
result, and the explanation in §6.5 is that the Reddit
\(r_{\text{true}}\) construction (political-lean magnitude
\(\times 0.25\), annotator divergence up to \(+0.3\), connection density
up to \(+0.1\)) blends contestedness with information density. The
community channel, not \(\hat{r}\), carries the contestedness signal in
the Stage 3 Scalable architecture. We read this as a measurement
decomposition that emerged from richer data, not a regression on the
underlying capability that Stage 2 demonstrated.

\subsubsection{5.10 Negative result: sharper supervised contrast
(v8b)}\label{negative-result-sharper-supervised-contrast-v8b}

Once v8a established that the encoder, freed from the prototype
bottleneck, organizes Reddit communities and external archetypes along a
structured trajectory manifold, a natural follow-up question was whether
\emph{more aggressive} supervised contrast would orthogonalize that
manifold further. The §5.9 archetype centroid off-diagonal cosine of
0.873 is well below v6/v7's 0.999 but still far from orthogonal, and the
within/between cosine ratio of 2.016, while doubled from v7's 1.006, is
also clearly improvable.

We trained v8b (10K steps, warm-started from v8a) with the community
supervised-contrastive loss weight raised from 2.0 to 4.0 and the
InfoNCE temperature lowered from 0.10 to 0.05. Every other architectural
and training choice was identical to v8a. Trainable parameter count is
unchanged at 14,560,579. The hypothesis: sharper contrastive pressure
should pull within-class cosines tighter and push between-class cosines
further, both improving Reddit retrieval and reducing archetype-centroid
alignment.

The result was a partial regression on every geometric metric.

{\def\LTcaptype{none} % do not increment counter
\begin{longtable}[]{@{}
  >{\raggedright\arraybackslash}p{(\linewidth - 6\tabcolsep) * \real{0.2500}}
  >{\raggedright\arraybackslash}p{(\linewidth - 6\tabcolsep) * \real{0.2500}}
  >{\raggedright\arraybackslash}p{(\linewidth - 6\tabcolsep) * \real{0.2500}}
  >{\raggedright\arraybackslash}p{(\linewidth - 6\tabcolsep) * \real{0.2500}}@{}}
\toprule\noalign{}
\begin{minipage}[b]{\linewidth}\raggedright
metric
\end{minipage} & \begin{minipage}[b]{\linewidth}\raggedright
v7
\end{minipage} & \begin{minipage}[b]{\linewidth}\raggedright
\textbf{v8a}
\end{minipage} & \begin{minipage}[b]{\linewidth}\raggedright
v8b
\end{minipage} \\
\midrule\noalign{}
\endhead
\bottomrule\noalign{}
\endlastfoot
val CE & 2.739 & 2.739 & 2.739 \\
Reddit recall@1 (35-cls) & 0.413 & \textbf{0.484} & 0.465 \\
within / between cosine ratio & 1.006 & \textbf{2.016} & 1.289 \\
archetype recall@1 (33-cls) & 0.149 & \textbf{0.230} & 0.214 \\
archetype centroid off-diag cosine & 0.999 & \textbf{0.873} & 0.945 \\
trajectory anisotropy (\(\lambda_{\max}/\lambda_{\min}\)) & 72 & 23,333
& \textbf{52,535} \\
regime ECE & 0.00091 & 0.00091 & \textbf{0.00070} \\
TruthfulQA mean\_r̂ AUROC & 0.578 & 0.577 & \textbf{0.579} \\
\end{longtable}
}

Cross-entropy and the BEN-side metrics (regime ECE, hallucination AUROC)
were preserved or marginally tightened. Every encoder-geometry metric
except trajectory anisotropy moved in the wrong direction. Reddit
within-class cosine pulled tighter (0.810 vs v8a's value), but
between-class cosine rose faster (0.628), so the ratio collapsed from
2.016 to 1.289. Archetype centroid off-diag cosine \emph{increased} from
0.873 back to 0.945, undoing roughly two-thirds of v8a's centroid
separation gain. Anisotropy more than doubled (52,535 vs 23,333),
indicating that the encoder collapsed a larger fraction of its variance
onto fewer principal directions.

The interpretation is that v8a's supcon weight 2.0 / temperature 0.10
was already near a sweet spot, and pushing harder reproduces a softer
version of the prototype-collapse failure one level up the architecture:
rather than collapsing 32 prototypes onto a handful of attractors, the
encoder collapses its 64-dimensional output onto a low-rank subspace
where a few directions carry most of the discriminative weight. The
contrastive objective, applied with too much pressure, optimizes a
degenerate solution that minimizes within-class spread by squashing the
entire embedding space.

We read v8b as a clean falsification of the ``sharper is better''
hypothesis. The continuous-trajectory architecture from v8a is the v8
generation's headline result; v8b documents the failure mode that bounds
it from above. Future work on the community head (Section 7) will not
pursue further increases in supcon weight or temperature sharpening on
this architecture, and will instead target either archetype-conditioned
direct supervision (Section 5.8 hypothesis (a)) or a fundamentally
different objective for orthogonalizing the trajectory manifold.

\subsubsection{5.12 Cross-backbone signal probe (Qwen 2.5-7B / Qwen 3-8B
/
Mistral-7B-v0.3)}\label{cross-backbone-signal-probe-qwen-2.5-7b-qwen-3-8b-mistral-7b-v0.3}

The single-backbone limitation flagged in §6.6 is partially addressable
without re-training the adapter on a new backbone. The first question to
ask is whether the discourse-community signal that v8a's CDH amplifies
is \emph{specific to Qwen 2.5}, or whether it is latent in other
7B-class causal LMs and merely needs to be read out. We answer with the
simplest possible probe: take raw mean-pooled hidden states from each
backbone on \texttt{data/val\_200.jsonl} (200 samples, 35 coarse Reddit
communities), then compute leave-one-out 1-nearest-neighbor recall@1 on
cosine similarity. No training, no labels at inference, no adapter.

{\def\LTcaptype{none} % do not increment counter
\begin{longtable}[]{@{}lrrr@{}}
\toprule\noalign{}
Backbone & best raw-hidden recall@1 & best layer & \(\times\) chance \\
\midrule\noalign{}
\endhead
\bottomrule\noalign{}
\endlastfoot
\texttt{Qwen/Qwen2.5-7B} & 0.260 & 16 & 9.1\(\times\) \\
\texttt{Qwen/Qwen3-8B} & 0.220 & 4 & 7.7\(\times\) \\
\texttt{mistralai/Mistral-7B-v0.3} & \textbf{0.295} & 16 &
\textbf{10.3\(\times\)} \\
\end{longtable}
}

For reference, chance on this 35-class probe is \(1/35 = 0.0286\), and
the v8a SRT-Adapter on \texttt{Qwen/Qwen2.5-7B} reaches recall@1
\(= 0.484\) (\(16.7\times\) chance) on the larger Reddit validation
split (§5.9). Three observations follow.

\emph{The signal is not Qwen-specific.} \texttt{Mistral-7B-v0.3} raw
hidden states actually carry slightly more discourse-community
information than \texttt{Qwen2.5-7B} raw hidden states do
(10.3\(\times\) vs 9.1\(\times\) chance on the same probe). The
community-discriminative subspace that the adapter's CDH learns to
amplify is present in both Qwen-family and Mistral-family backbones at
comparable strength. Whatever the adapter is reading off the residual
stream, the \emph{raw substrate} it depends on is a property of
contemporary 7B-class causal LMs, not of one specific tokenizer or
training recipe.

\emph{The adapter's contribution is a \textasciitilde{}\(1.8\times\)
amplification of that substrate.} v8a on Qwen 2.5-7B delivers
\(16.7\times\) chance, which is roughly \(1.8\times\) the strongest
raw-hidden-state baseline in the table (Mistral, \(10.3\times\)) and
\(1.8\times\) the same backbone's own raw signal at the best layer
(\(9.1\times\)). The adapter is doing real discriminative work on top of
what the frozen backbone already encodes; the residual stream alone gets
you most of the way to chance-level retrieval but does not reach the v8a
number on its own.

\emph{The architecture should transfer.} The CDH consumes a single
hidden-state tensor at \texttt{community\_layer\_idx}; nothing in the
adapter is Qwen-tokenizer-specific. Re-training the v8a recipe on
Mistral-7B-v0.3 is a plausible future-work item; the raw-substrate
result above is the strongest zero-cost evidence we currently have for
the cross-backbone hypothesis. The full reproduction script is
\texttt{scripts/cross\_backbone\_probe.py}, JSON results are at
\texttt{artifacts/cross\_backbone.json} and the two extra-backbone files
(\texttt{cross\_backbone\_extra.json},
\texttt{cross\_backbone\_mistral.json}), and the rendered visualization
is Fig. 8 of the public demo card.

\subsubsection{5.13 Interiority study: post-release diagnostics on
v8a}\label{interiority-study-post-release-diagnostics-on-v8a}

This subsection summarizes a sequence of post-release probes on the v8a
checkpoint that dissect \emph{how} the adapter encodes semiotic regime,
with full numbers and reproduction commands in
\href{https://github.com/space-bacon/SRT/blob/main/docs/INTERIORITY_V1_FINDINGS.md}{\texttt{docs/INTERIORITY\_V1\_FINDINGS.md}}.
The probe battery is 11 regimes (\texttt{metaphor},
\texttt{counterfactual}, \texttt{negation\_modality}, \texttt{irony},
\texttt{literal}, \texttt{refusal\_bait}, \texttt{lyric},
\texttt{deixis}, \texttt{self\_reference}, \texttt{quoted\_speech},
\texttt{code}) \(\times\) 25 prompts. The community vector throughout is
the 64-D CDH output (\texttt{use\_prototypes=False}).

\emph{Layer-as-organ map.} Per-regime z-score winners across nine
adapter channels show channel specialization: \texttt{code} peaks on
\(\hat{r}\), regime entropy, and \(\text{div}\,L21\); \texttt{metaphor}
peaks on \(\text{inj}\,L14\), \(\text{inj}\,L21\), and the chain
residual; \texttt{literal} peaks on \(P(\text{super})\). The four-module
decomposition acquires functionally distinct channels under v8a's
continuous architecture.

\emph{Community 1-NN.} Leave-one-out 1-NN top-1 accuracy on the
11-regime probe battery is \(0.70\), with bootstrap mean \(0.77\) and
CI95 \([0.72, 0.82]\) over 2000 iterations. This is \(7.7\times\) chance
on a 1/11 task. The top centroid pair
\texttt{code}\(\leftrightarrow\)\texttt{literal} (Euclidean
\(d = 1.07\)) stays the most-separated pair in \(85.5\%\) of bootstrap
iterations; the bottom pair
\texttt{quoted\_speech}\(\leftrightarrow\)\texttt{self\_reference}
(\(d = 0.39\)) stays last in \(83.0\%\). The point estimate of community
separation ratio is \(0.85\), which we treat as a lower bound: the
bootstrap CI95 \([0.86, 0.98]\) is upward-biased by within-cluster
duplicate sampling.

\emph{BOS-register length scaling.} A per-token trajectory probe across
\(T \in \{16, 49, 79, 129\}\) token windows (4 length bins \(\times\) 5
prompts \(\times\) 11 regimes) shows that the adapter's write channels
(\(\text{inj}\,L14\), \(\text{inj}\,L21\), \(\text{div}\,L7/L14/L21\),
chain residual) form a \emph{fixed-size BOS register}: BOS-token
amplitude on those channels is stable across an \(8\times\) length
sweep, staying within \(\pm 2.5\%\) of the \(T \approx 16\) baseline
(peak-to-peak \(\leq 4.8\%\)). v8a learns to compress each prompt's
regime structure into a single fixed-amplitude write at the BOS slot;
mid-prompt activity on inject channels is small.

\emph{BOS-sink ablation, per channel \(\times\) regime.} The above is
consistent with two readings: either the adapter encodes its global
regime summary at BOS and only at BOS, or BOS is a sink that contributes
to but does not solely determine the regime ranking. We disambiguate by
computing, for each channel \(c\) and regime \(g\), the z-score of
\(g\)'s mean activation against the other ten regimes twice: once with
BOS included in the pooling window (\(z^{\text{with}}_{c,g}\)), once
with BOS excluded (\(z^{\text{no}}_{c,g}\)), and reporting
\(\Delta |z|_{c,g} = |z^{\text{with}}_{c,g}| - |z^{\text{no}}_{c,g}|\).
On \(\hat{r}\) and on regime entropy, \(|\Delta|z|| \leq 0.27\) for
every regime: the regime-classifier scalars read essentially the same
regime structure with or without the BOS slot, so the regime signature
is \emph{content-borne}. On \(\text{inj}\,L14\) the deltas are large and
signed: \texttt{metaphor} \(\Delta = -0.64\) (regime signal is
BOS-localized; removing BOS \emph{strengthens} the relative z-score of
the content-token mean), \texttt{literal} \(\Delta = +0.98\) (regime
signal is \emph{anti}-localized to BOS; removing BOS dilutes the regime
signature). The chain-residual channel is intermediate (max \(|\Delta|\)
\(= 0.58\) on \texttt{quoted\_speech}). The honest synthesis is that the
inject channels are functioning as register slots that hold a global
per-prompt summary, while the regime classifier reads regime structure
off content tokens directly. A simpler aggregate measure, regime-rank
preservation under BOS exclusion, holds in 9/9 channels with mean
Spearman \(\rho = +0.74\), which is consistent with but coarser than the
per-channel per-regime breakdown above.

\emph{Within-prompt trajectories.} Plotting per-position curves for the
regime-classifier scalars (\(\hat{r}\), regime entropy) across 19
content position bins (BOS dropped) per regime shows that the temporal
structure hidden by mean pooling is informative. \(\hat{r}\) separates
regimes by sustained amplitude: \texttt{code} sits at
\(\hat{r} \geq 0.85\) across the full window;
\texttt{negation\_modality}, \texttt{deixis}, and \texttt{refusal\_bait}
sit at \(\hat{r} \leq 0.40\); the other seven regimes occupy a middle
band \(0.45\)--\(0.75\). Regime entropy puts \texttt{code} in a
sustained \(0.18\)--\(0.32\) band while every other regime stays at
\(\leq 0.05\) across all 19 content positions. The classifier is
consistently \emph{uncertain} about the regime label whenever the prompt
is \texttt{code} and consistently confident on every other regime, and
that uncertainty is steady-state across positions, not a transient.
Mean-pooled scalars collapse this to a single number per prompt and lose
the structure.

\emph{Architectural reading of the interiority study.} v8a learned to
compress each prompt into a fixed-size \emph{regime fingerprint} written
into the BOS register on inject channels, while the regime classifier
itself operates as a content-distributed channel that reads regime
evidence off non-BOS tokens. This is consistent with §6.3: the
inject-back arm carries no measurable signal in v8a, so the adapter's
content-routing capacity is under-used. v9 (Section 7) targets this
directly with a coverage-loss term that penalizes register
concentration. The interiority probes do not change v8a's headline
numbers; they explain \emph{where} in the architecture those numbers
come from. Visualizations (per-channel z-score heatmaps, per-regime
within-prompt trajectories, BOS-ablation \(\Delta |z|\) matrix) appear
as Figs. 10 and 11 of the public demo card.

\begin{center}\rule{0.5\linewidth}{0.5pt}\end{center}

\subsection{6. Discussion}\label{discussion}

\subsubsection{6.1 Semiotic Structure in Frozen
Representations}\label{semiotic-structure-in-frozen-representations}

The adapter architecture embodies a specific theoretical claim: that the
semiotic structure of discourse, comprising community-conditioned
interpretations, divergence patterns, and bifurcation dynamics, is
already encoded in the hidden states of pretrained language models. The
claim follows necessarily from the fact that these models were trained
on text produced by communities with divergent interpretive norms. What
the adapter adds is not new information but new \emph{readout
apparatus}: projections, attention mechanisms, and recurrence that
disentangle the semiotic structure already present.

This is analogous to the relationship between a microscope and the
structures it reveals. The adapter does not create bifurcation dynamics
in text. It provides the lenses through which dynamics that were always
present become visible and measurable.

\subsubsection{6.2 Community Discovery Without
Labels}\label{community-discovery-without-labels}

A significant departure from the original SRT is the replacement of
supervised community embeddings with unsupervised prototype-based
clustering. The original architecture required explicit community IDs at
both training and inference time, limiting deployment to domains with
known community structure. The adapter's community head learns to
partition discourse space from backbone hidden states alone, discovering
whatever grouping structure best serves the downstream semiotic losses.

This is more faithful to Peirce's framework, in which communities of
interpretation are not given \emph{a priori} but emerge through shared
interpretive practice. The prototypes are pulled apart by the semiotic
losses: if assigning text to different communities helps the model
predict divergence better, it will learn to separate them. Community
structure is discovered, not imposed.

\subsubsection{6.3 The Injection Pathway: Observation
vs.~Intervention}\label{the-injection-pathway-observation-vs.-intervention}

The RRM's injection mechanism creates a feedback loop between semiotic
observation and language generation. This is the architectural
instantiation of metapragmatic awareness: the model's observation of
divergence changes the hidden states that produce subsequent text. The
injection is deliberately small (scale factor \(\alpha = 0.1\),
zero-initialized projection, sigmoid gating), reflecting a conservative
design philosophy: the adapter should primarily \emph{observe} semiotic
dynamics. Active intervention, that is, generation that \emph{responds}
to detected bifurcation, is an advanced capability that requires careful
validation before scaling.

The CE loss provides a natural safety valve. Since gradients from CE
flow through the injection pathway, the model is penalized if injections
degrade language modeling quality. This creates an automatic pressure
toward injections that are either helpful or neutral, never harmful.

The empirical situation through v8b is that this pressure has resolved
on the \emph{neutral} side of the helpful/neutral boundary: ablating the
inject-back arm at evaluation produces no measurable downstream change
on CE, on the regime classifier, or on the hallucination probes
(Sections 2.5, 5.7, 5.9). The safety valve held, but the channel through
which the meta-state would modulate generation has not learned to carry
signal. The interiority study in §5.13 sharpens this picture: the inject
channels \emph{do} carry signal, but it is concentrated in a fixed-size
BOS register (BOS amplitude is stable across an \(8\times\) length
sweep, peak-to-peak \(\leq 4.8\%\)), with the per-channel \(\times\)
per-regime BOS-ablation analysis showing that some regimes
(e.g.~\texttt{metaphor}) localize their inject signature at BOS and
others (e.g.~\texttt{literal}) are anti-localized to BOS. The arm has
learned to \emph{summarize} but not to \emph{modulate}. v9's
coverage-loss design (Section 7, Appendix A version history) is directed
at exactly this: penalize register concentration and credit distributed
mid-prompt activity, so that the inject channel has gradient pressure to
behave as a per-token modulator rather than a per-prompt summary slot.

This null is not new. Stage 3 Phase 1 ran 105 training rounds (R21
through R105 on TinyLlama-1.1B, documented in Lancaster, 2026a) with the
explicit goal of activating the inject-back arm, including a remediation
campaign covering loss-weight sweeps, BEN architecture overhauls, FiLM
scale schedules, and gradient-isolation experiments. The arm did not
activate on TinyLlama and has not activated through v8b on Qwen. Two
readings are consistent with the data: (1) FiLM injection requires
different scaffolding when the backbone is frozen, suggesting
alternative injection mechanisms (cross-attention from RRM meta-state
into selected backbone layers, low-rank residual-stream modulation
conditioned on \(\hat{r}\), learned gating bypassed by a meta-state
classifier) are the appropriate v9-onward design target; (2) the mutual
information carried by RRM meta-state about downstream gradients is
fundamentally small when the backbone is frozen (per the
information-theoretic bound in Section 2.6), placing a low ceiling on
the intervention arm's possible effectiveness. We do not have the
experimental record to distinguish these readings yet. v9 onward targets
the first hypothesis directly. The current adapter should be
characterized as a self-organizing observation channel over a frozen
backbone, not yet as a closed circular-causal system.

\subsubsection{6.4 Relation to the Pitchfork
Model}\label{relation-to-the-pitchfork-model}

BEN's \(\hat{r}\) estimate is the primary output of the entire system.
It provides a per-token, continuous measure of semiotic stability that
maps directly onto the control parameter of the pitchfork bifurcation:

\begin{itemize}
\tightlist
\item
  \(\hat{r} < 0\): The sign is in the subcritical regime. Shared meaning
  is stable. Perturbations decay.
\item
  \(\hat{r} \approx 0\): The sign is near-critical. Small changes in
  context or community could tip it.
\item
  \(\hat{r} > 0\): The sign has bifurcated. Meaning has split into
  community-specific attractors.
\end{itemize}

This is not a classifier applied after the fact. \(\hat{r}\) is
estimated from the accumulated meta-state of the RRM, which tracks how
divergence has evolved through the backbone's processing hierarchy. It
is a real-time structural estimate, not a post-hoc label.

\subsubsection{\texorpdfstring{6.5 What \(\hat{r}\) actually
measures}{6.5 What \textbackslash hat\{r\} actually measures}}\label{what-hatr-actually-measures}

The context-conditional probe in §5.6 returned a null result at the
target token and a \emph{negative} result over the full passage. This is
informative. Earlier drafts of this paper described \(\hat{r}\) as a
contestedness detector, that is, as a per-token estimate of how much a
sign is being fought over in its current discourse register. The data
does not support that interpretation.

What \(\hat{r}\) actually appears to measure is information density
combined with the specific reflexivity components encoded in
\(r_{\text{true}}\) (political-lean magnitude, annotator divergence,
connection density). Fact-dense prose (dates, numbers, named entities,
citations) drives \(\hat{r}\) up because those are the positions where
\(r_{\text{true}}\) is highest in the training distribution. Rhetorical
or formulaic charged language is, on average, \emph{less} lexically
diverse than dense factual writing, so its mean \(\hat{r}\) is lower.

The contestedness signal is in the \emph{community head}, not in
\(\hat{r}\). The counterfactual-decoding result (§5.3) and the per-topic
community shifts in §5.6 both demonstrate this. The v6 divergence-SupCon
loss was intended to sharpen the metapragmatic channel further; in
practice it improved community recall@1 (§5.7) but did not produce the
expected community-conditional separation of MAH divergence trajectories
on contested topics, and v8a's prototype-bottleneck removal (§5.9)
turned out to be the more consequential change for both channels.

This revision of the architectural story is cleaner, not weaker: the
model has two distinct outputs that measure two distinct things, and the
data tells us which is which.

\subsubsection{6.6 Limitations}\label{limitations}

\begin{enumerate}
\def\labelenumi{\arabic{enumi}.}
\item
  \textbf{No modulation at inference.} The current architecture
  estimates \(\hat{r}\) but does not use it to modulate generation.
  Future work will explore \(\lambda\)-controlled modes where detected
  bifurcation triggers bridge-generation strategies.
\item
  \textbf{Simplified regime model.} The binary subcritical/supercritical
  classification omits the near-critical regime, which is arguably the
  most important for practical applications (early warning of emerging
  bifurcation). The three-class model from the original SRT will be
  restored once binary classification is validated.
\item
  \textbf{Reddit-only data.} Training on Reddit discourse may not
  generalize to other domains (news media, academic text, legal
  documents). The TruthfulQA hallucination probe (§5.4) is a partial
  cross-domain transfer signal but mean \(\hat{r}\) AUROC of 0.573 is
  well below the 0.7 production threshold.
\item
  \textbf{No human evaluation.} All supervision comes from computed
  \(r_{\text{true}}\) labels. Ecological validity, that is, whether
  \(\hat{r}\) tracks what human annotators perceive as meaning
  contestation, has not been tested. The §5.6 negative result is the
  strongest current evidence that \(r_{\text{true}}\) as currently
  constructed is not a clean proxy for contestedness.
\item
  \textbf{Single backbone (partially answered).} Adapter-trained results
  are reported only for Qwen 2.5-7B. The cross-backbone raw-hidden probe
  in §5.12 shows the discourse-community signal that v8a amplifies is
  present at comparable strength in \texttt{Qwen/Qwen3-8B}
  (\(7.7\times\) chance on raw hidden states) and
  \texttt{mistralai/Mistral-7B-v0.3} (\(10.3\times\) chance, slightly
  stronger than Qwen 2.5-7B's own \(9.1\times\)). This is consistent
  with the architecture being backbone-agnostic and with the adapter's
  \(\sim 1.8\times\) amplification factor transferring, but the actual
  adapter-on-Mistral re-train has not yet been executed. The full
  backbone-agnostic claim still requires that re-training, which is a
  v9-onward work item.
\end{enumerate}

\subsubsection{6.7 Connections to Emergent Perspective Diversity in
Reasoning
Models}\label{connections-to-emergent-perspective-diversity-in-reasoning-models}

A separate line of recent work has documented that language models
trained with reasoning-style reinforcement spontaneously develop
heterogeneous internal features that mechanistic interpretability
methods can read out as something like distinct personalities, areas of
expertise, or stances, and that these features enter into structured
conflict and reconciliation during chain-of-thought generation (Evans,
in preparation; cf.~Foster, Rzhetsky, \& Evans, 2015 for the
cross-corpus precedent). The relationship between that finding and the
present work is structural rather than methodological. Evans's program
reveals what emerges \emph{spontaneously} under reasoning supervision in
an architecturally-undifferentiated backbone; the SRT-Adapter provides
explicit architectural channels (the four Peircean subspaces of the
semiotic embedding layer in the full SRT, and the community / divergence
/ bifurcation subspaces of the present adapter) into which similar
emergent structure can organize. The two approaches are complementary in
a specific sense: probing the unstructured backbone reveals the
existence of perspective-like features without a vocabulary for what
they are; the adapter's explicit decomposition supplies a vocabulary,
grounded in Peircean semiotics and linguistic anthropology, but only
weakly constrains what fills the slots.

This suggests a research question that neither approach can resolve
alone. Reasoning models that develop internal perspective diversity may
be performing genuine semiotic work, namely navigating meaning
divergence across implicit interpretive communities, or they may be
performing something closer to cognitive brainstorming or rhetorical
variation within a single interpretive framework. The SRT-Adapter's
metapragmatic attention head is designed to detect the former and would,
in principle, register low divergence across the latter. Cross-applying
mechanistic interpretability tools to adapter-equipped reasoning models,
and cross-applying the adapter's divergence and bifurcation readouts to
backbones probed for emergent features, is the form of collaboration
this paper points toward as the next step.

The von Foerster framing (Section 2.5) makes the same point in
cybernetic terms: spontaneous perspective diversity in reasoning models
is self-organization without an explicit semiotic substrate; the adapter
provides the substrate without yet exhibiting the fully closed
circular-causal loop. The two findings, taken together, suggest that the
substrate and the emergent dynamics may be separable engineering
targets, and that the productive question is what becomes possible when
both are present.

\begin{center}\rule{0.5\linewidth}{0.5pt}\end{center}

\subsection{7. Conclusion}\label{conclusion}

The SRT-Adapter demonstrates that semiotic awareness can be added to any
frozen language model as a lightweight, modular capability. By tapping
hidden states rather than rebuilding the backbone, the architecture
preserves pretrained language modeling quality (CE = 2.63 vs.~unadapted
2.71 on the same val data) while introducing structured outputs that
make the semiotic dynamics of text visible and measurable.

The v5 generation establishes the basic capability set on five
independent probes:

\begin{enumerate}
\def\labelenumi{\arabic{enumi}.}
\tightlist
\item
  \textbf{CE preservation} (§5.1): the injection pathway is mildly
  helpful, not harmful.
\item
  \textbf{Community geometry} (§5.2): recall@1 of 0.36 on a 35-class
  unsupervised retrieval task (\(12.6\times\) random), made possible by
  SupCon on the encoder's pre-mixing output.
\item
  \textbf{Counterfactual community decoding} (§5.3): forcing the
  community vector at decode time produces zero disagreement on factual
  prompts and near-total disagreement on contested topics, demonstrating
  the community vector behaves as a discourse prior.
\item
  \textbf{Hallucination signal} (§5.4): all four SRT-native channels
  lean in the predicted direction on TruthfulQA without truthfulness
  supervision; mean \(\hat{r}\) AUROC = 0.573.
\item
  \textbf{Regime calibration} (§5.5): ECE = \(9 \times 10^{-4}\) and
  AUROC = 0.99 on 351K tokens, unblocking downstream probabilistic use.
\end{enumerate}

The negative result on context-conditional \(\hat{r}\) (§5.6) sharpened
the architectural story: \(\hat{r}\) measures information density and
the components of \(r_{\text{true}}\), while the community head measures
discourse register. Conflating these in earlier drafts was a theoretical
error the data corrected.

The v6--v8 generations then reframe the design. v6 and v7 (§5.7) extend
SupCon from the community channel to the metapragmatic divergence
channel and add ListNet ranking on \(\hat{r}\), giving incremental gains
on community recall@1 (0.360 → 0.413) and calibration (ECE → 0.0006).
The cross-corpus convergence probe (§5.8) shows the prototype tensors
barely move during training and that 33 externally-curated archetypes
collapse onto roughly four functional macro-clusters, reproducing at
single-backbone scale the macro-attractor pattern Foster, Rzhetsky, \&
Evans (2015) document for scientific subdisciplines. v8a (§5.9) is the
headline architectural result: removing the 32-prototype mixing layer
entirely leaves CE unchanged (\(\Delta = +0.0001\) nats) while raising
Reddit recall@1 from 0.413 to 0.484, raising archetype recall@1 to
\(7.6\times\) chance, nearly doubling the within/between cosine ratio,
and expanding trajectory anisotropy by \(\sim 325\times\). v8b (§5.10)
falsifies the ``sharper-supcon'' hypothesis on the continuous-encoder
architecture, bounding the v8 design from above. The encoder, not the
prototype basis, was doing the discriminative work; the prototype layer
was discarding it through a saturated soft-argmax.

What the v5--v8 arc has demonstrated is a \emph{self-organizing
observation channel over a frozen backbone} (Sections 2.5--2.6,
5.1--5.10). What it has not yet demonstrated is a \emph{closed
circular-causal loop} in which the meta-state's observation modifies
generation in a measurable way: ablating the inject-back arm produces no
downstream change on CE, calibration, or hallucination probes through
v8b (Sections 2.5, 6.3). Reading this through the information-bound
framing of Section 2.6, the simplest hypothesis is that the mutual
information the meta-state currently carries about the downstream loss
is small, and that pushing \(Q_0\) above the bifurcation threshold of
the inject-back arm requires either a larger meta-state, a different
injection geometry, or direct supervision on the closed-loop behavior.
v9 onward is the design target for that work, together with
archetype-conditioned direct supervision (§5.8 hypothesis (a)) for
resolving sub-macro-cluster archetype structure. Combined-feature
hallucination probes, cross-domain transfer evaluation, and human
ecological-validity studies remain the principal open empirical
questions.

\begin{center}\rule{0.5\linewidth}{0.5pt}\end{center}

\subsection{References}\label{references}

Agha, A. (2003). The social life of cultural value. \emph{Language \&
Communication}, 23(3--4), 231--273.

Anderson, M. (2014). Mathematical modeling of catastrophic change in
cultural systems. In M. Anderson (Ed.), \emph{Cultural shaping of
violence: Victimization, escalation, response} (selected chapters).
Purdue University Press.

Bail, C. A., et al.~(2018). Exposure to opposing views on social media
can increase political polarization. \emph{Proceedings of the National
Academy of Sciences}, 115(37), 9216--9221.

Bennett, C. H. (1982). The thermodynamics of computation: A review.
\emph{International Journal of Theoretical Physics}, 21(12), 905--940.

Deely, J. (2014). The suprasubjective in semiotic relations. \emph{The
American Journal of Semiotics}, 30(3--4), 165--182.

Durst-Andersen, P. (2011). \emph{Linguistic supertypes: A
cognitive-semiotic theory of human communication}. De Gruyter Mouton.

Evans, J. A. (2010). Industry induces academic science to know less
about more. \emph{American Journal of Sociology}, 116(2), 389--452.

Foster, J. G., Rzhetsky, A., \& Evans, J. A. (2015). Tradition and
innovation in scientists' research strategies. \emph{American
Sociological Review}, 80(5), 875--908.

Irvine, J. T., \& Gal, S. (2000). Language ideology and linguistic
differentiation. In P. V. Kroskrity (Ed.), \emph{Regimes of language}
(pp.~35--83). SAR Press.

Kockelman, P. (2017). \emph{The art of interpretation in the age of
computation}. Oxford University Press.

Kockelman, P. (2024). \emph{Last words: A theory of everything that
matters}. University of Chicago Press.

Kockelman, P. (2025). \emph{Semiotic agency in digital environments}.
Manuscript.

Lancaster, J. B. (2025). The treachery of signs: Semiotic mediation,
pitchfork bifurcation, and political polarization in algorithmically
curated societies. SSRN. https://papers.ssrn.com/abstract=5987495

Lancaster, J. B. (2026a). Semiotic-reflexive language model training:
Bridging interpretive bifurcations through metapragmatic chain
architectures and embodied grounding. SSRN.
https://papers.ssrn.com/abstract=6349978

Lancaster, J. B. (2026b). Prenatal origins of cross-modal iconic
correspondence: A semiotic analysis.

Lancaster, J. B. (2026c). Reddit Discourse Corpus: A multi-community
dataset for semiotic analysis.

Landauer, R. (1961). Irreversibility and heat generation in the
computing process. \emph{IBM Journal of Research and Development}, 5(3),
183--191.

Latour, B. (1996). On interobjectivity. \emph{Mind, Culture, and
Activity}, 3(4), 228--245.

Leighton, M. P. (2026). Will a large complex system be a Maxwell demon?
\emph{arXiv preprint} arXiv:2603.03248.

Mangalam, M. (2025). Against the Bayesian brain. \emph{Behavioral and
Brain Sciences} (forthcoming).

Maturana, H. R., \& Varela, F. J. (1980). \emph{Autopoiesis and
cognition: The realization of the living}. D. Reidel.

Parrondo, J. M. R., Horowitz, J. M., \& Sagawa, T. (2015).
Thermodynamics of information. \emph{Nature Physics}, 11(2), 131--139.

Peirce, C. S. (1931--1958). \emph{Collected papers of Charles Sanders
Peirce} (Vols. 1--8). C. Hartshorne, P. Weiss, \& A. Burks (Eds.).
Harvard University Press.

Radford, A., et al.~(2021). Learning transferable visual models from
natural language supervision. In \emph{ICML 2021}.

Ramachandran, V. S., \& Hubbard, E. M. (2001). Synaesthesia: a window
into perception, thought and language. \emph{Journal of Consciousness
Studies}, 8(12), 3--34.

Silverstein, M. (1993). Metapragmatic discourse and metapragmatic
function. In J. A. Lucy (Ed.), \emph{Reflexive language} (pp.~33--58).
Cambridge University Press.

Silverstein, M. (2003). Indexical order and the dialectics of
sociolinguistic life. \emph{Language \& Communication}, 23(3--4),
193--229.

VanSaders, B., Fruchart, M., \& Vitelli, V. (2026). Measurement-induced
phase transitions in informational active matter. \emph{PNAS Nexus},
pgag077. https://doi.org/10.1093/pnasnexus/pgag077

Versace, E., et al.~(2023). Cross-modal correspondences between auditory
and visual features in domestic chicks. \emph{Animal Cognition}, 26,
1021--1030.

von Foerster, H. (1981). \emph{Observing systems}. Intersystems
Publications.

von Foerster, H. (2003). \emph{Understanding understanding: Essays on
cybernetics and cognition}. Springer.

Wildgen, W. (1982). \emph{Catastrophe-theoretic semantics: An
elaboration and application of René Thom's theory}. John Benjamins.

\begin{center}\rule{0.5\linewidth}{0.5pt}\end{center}

\subsection{Appendix A: Configuration Defaults
(v8a)}\label{appendix-a-configuration-defaults-v8a}

\begin{Shaded}
\begin{Highlighting}[]
\NormalTok{SRTConfig(}
\NormalTok{    backbone\_id    }\OperatorTok{=} \StringTok{"Qwen/Qwen2.5{-}7B"}\NormalTok{,}
\NormalTok{    backbone\_dtype }\OperatorTok{=} \StringTok{"bfloat16"}\NormalTok{,}
\NormalTok{    mah }\OperatorTok{=}\NormalTok{ MAHConfig(d\_sub}\OperatorTok{=}\DecValTok{512}\NormalTok{, d\_divergence}\OperatorTok{=}\DecValTok{256}\NormalTok{, num\_heads}\OperatorTok{=}\DecValTok{4}\NormalTok{, dropout}\OperatorTok{=}\FloatTok{0.1}\NormalTok{),}
\NormalTok{    rrm }\OperatorTok{=}\NormalTok{ RRMConfig(d\_meta}\OperatorTok{=}\DecValTok{512}\NormalTok{, inject\_scale}\OperatorTok{=}\FloatTok{1.0}\NormalTok{),  }\CommentTok{\# FiLM since v4}
\NormalTok{    ben }\OperatorTok{=}\NormalTok{ BENConfig(d\_hidden}\OperatorTok{=}\DecValTok{256}\NormalTok{),                  }\CommentTok{\# tanh removed in v4}
\NormalTok{    community }\OperatorTok{=}\NormalTok{ CommunityConfig(}
\NormalTok{        num\_prototypes}\OperatorTok{=}\DecValTok{32}\NormalTok{, d\_community}\OperatorTok{=}\DecValTok{64}\NormalTok{, temperature}\OperatorTok{=}\FloatTok{1.0}\NormalTok{,}
\NormalTok{        use\_prototypes}\OperatorTok{=}\VariableTok{False}\NormalTok{,                       }\CommentTok{\# v8a: encoder output IS the community vector}
\NormalTok{    ),}
\NormalTok{    loss }\OperatorTok{=}\NormalTok{ LossConfig(}
\NormalTok{        ce\_weight}\OperatorTok{=}\FloatTok{1.0}\NormalTok{, chain\_weight}\OperatorTok{=}\FloatTok{0.5}\NormalTok{, bif\_weight}\OperatorTok{=}\FloatTok{1.0}\NormalTok{,}
\NormalTok{        regime\_weight}\OperatorTok{=}\FloatTok{5.0}\NormalTok{, div\_alive\_weight}\OperatorTok{=}\FloatTok{0.1}\NormalTok{,}
\NormalTok{        inject\_reg\_weight}\OperatorTok{=}\FloatTok{0.0}\NormalTok{, inject\_target\_norm}\OperatorTok{=}\FloatTok{1.0}\NormalTok{,    }\CommentTok{\# v4: dropped (FiLM init handles it)}
\NormalTok{        community\_entropy\_weight}\OperatorTok{=}\FloatTok{0.01}\NormalTok{,}
        \CommentTok{\# v5}
\NormalTok{        community\_supcon\_weight}\OperatorTok{=}\FloatTok{2.0}\NormalTok{,}
\NormalTok{        community\_supcon\_temperature}\OperatorTok{=}\FloatTok{0.1}\NormalTok{,}
        \CommentTok{\# v6}
\NormalTok{        divergence\_supcon\_weight}\OperatorTok{=}\FloatTok{0.3}\NormalTok{,                     }\CommentTok{\# v7: dropped from 1.0 to recover §5.3}
\NormalTok{        divergence\_supcon\_temperature}\OperatorTok{=}\FloatTok{0.1}\NormalTok{,}
\NormalTok{        listnet\_weight}\OperatorTok{=}\FloatTok{0.5}\NormalTok{,}
\NormalTok{        listnet\_temperature}\OperatorTok{=}\FloatTok{1.0}\NormalTok{,}
\NormalTok{        chain\_residual\_aux\_weight}\OperatorTok{=}\FloatTok{0.05}\NormalTok{,}
\NormalTok{        chain\_residual\_aux\_target}\OperatorTok{=}\FloatTok{0.5}\NormalTok{,}
\NormalTok{    ),}
\NormalTok{)}
\end{Highlighting}
\end{Shaded}

\subsubsection{Version history}\label{version-history}

{\def\LTcaptype{none} % do not increment counter
\begin{longtable}[]{@{}
  >{\raggedright\arraybackslash}p{(\linewidth - 4\tabcolsep) * \real{0.3333}}
  >{\raggedright\arraybackslash}p{(\linewidth - 4\tabcolsep) * \real{0.3333}}
  >{\raggedright\arraybackslash}p{(\linewidth - 4\tabcolsep) * \real{0.3333}}@{}}
\toprule\noalign{}
\begin{minipage}[b]{\linewidth}\raggedright
Version
\end{minipage} & \begin{minipage}[b]{\linewidth}\raggedright
Headline change
\end{minipage} & \begin{minipage}[b]{\linewidth}\raggedright
Result
\end{minipage} \\
\midrule\noalign{}
\endhead
\bottomrule\noalign{}
\endlastfoot
v1 & Initial architecture, \(\tanh\) on \(\hat{r}\), L2 on injections &
\(\hat{r}\) saturated at \(\pm 1\), injections at norm \(\approx 7\) \\
v2 & Diagnostic instrumentation added & Confirmed both pathologies; CE
healthy \\
v3 & Target-norm injection penalty \((\|\text{inj}\| - 1)^2\) &
Injections recovered; community prototypes still collapsed (recall@1
\(= 0.05\)) \\
v4 & \(\tanh\) removed from BEN; RRM linear-gated → FiLM; first SupCon
attempt on \texttt{vector} (failed) & \(\hat{r}\) tail recovered; SupCon
flatlined at \(\log(B-1)\) \\
v5 & SupCon switched to \texttt{encoded} (pre-mixing); weight
\(0.5 \to 2.0\); warm-restart of community head & recall@1 \(= 0.36\),
ECE \(= 9 \times 10^{-4}\), counterfactual-decode contested/factual
split \\
v6 & Add divergence-SupCon (\(\lambda = 1.0\)), ListNet on \(\hat{r}\),
chain-residual auxiliary floor & Reddit recall@1 \(0.36 \to 0.41\), ECE
\(\to 0.0006\); §5.3 counterfactual decode regressed \\
v7 & Reduce divergence-SupCon to \(\lambda = 0.3\) & Reddit recall@1
\(\to 0.413\), hallu AUROC \(\to 0.5785\), decode signal recovered \\
v8a & Drop the 32-prototype mixing layer
(\texttt{use\_prototypes=False}), encoder output = community vector & CE
\(\Delta = +0.0001\); Reddit recall@1 \(\to 0.484\); archetype recall@1
\(7.6\times\) chance; within/between cosine ratio \(\to 2.016\);
trajectory anisotropy \(\times 325\) \\
v8b & Sharpen community-SupCon (\(\lambda\colon 2.0 \to 4.0\),
\(\tau\colon 0.10 \to 0.05\)) on v8a base & Partial regression on every
encoder-geometry metric except anisotropy; falsifies ``sharper is
better'' on this architecture \\
v8a interiority (post-release) & Eight diagnostic probes on the v8a
checkpoint (§5.13): layer-as-organ map, community 1-NN, BOS-register
length scaling, BOS-sink ablation per channel \(\times\) regime,
within-prompt trajectories, regime calibration replication,
cross-backbone raw-hidden probe (§5.12) & Confirms v8a inject channels
function as a fixed-size BOS register (peak-to-peak \(\leq 4.8\%\)
across \(8\times\) length sweep); regime classifier reads content-borne
signal; calibration ECE replicates at \(9.1\times 10^{-4}\) on 351K
tokens; community 1-NN at \(7.7\times\) chance; cross-backbone signal
latent in Qwen 3-8B and Mistral-7B at comparable strength \\
v9 (in progress) & Closed-loop training target for the inject-back arm
(Sections 2.5, 6.3); coverage loss penalizing register concentration
(§5.13, §6.3); archetype-conditioned direct supervision (§5.8 hypothesis
(a)); adapter re-train on Mistral-7B-v0.3 (§6.6) & TBD \\
\end{longtable}
}

\subsection{Appendix B: Layer Index
Auto-Computation}\label{appendix-b-layer-index-auto-computation}

Given backbone depth \(L\): - MAH hook layers:
\([\lfloor L/4 \rfloor, \, \lfloor L/2 \rfloor, \, \lfloor 3L/4 \rfloor]\)
- RRM injection layers: MAH layers 2 and 3 (skip first to let meta-state
accumulate) - Community discovery layer:
\(\max(1, \lfloor L/7 \rfloor)\)

For Qwen 2.5-7B (\(L = 28\)): MAH @ {[}7, 14, 21{]}, inject @ {[}14,
21{]}, community @ 4.

\end{document}
